
\label{ch:translational}


\begin{ataglance}[How to read this chapter]
This chapter does one thing. It takes a construct built on a general ageing
cohort and asks what would have to be true for it to work in haemophilia. No
haemophilia data are analysed, so what follows is a set of commitments made
before the data exist.

The order is a chain of prerequisites. \S\ref{sec:why_haemophilia} says why
haemophilia is the test case worth attempting first. \S\ref{sec:hemo_mapping}
replaces the general indicators with clinical ones and states, in
Derivation~\ref{der:indicator_substitution}, the price of doing so.
\S\ref{sec:two_stage} checks that swapping indicators has not quietly turned
translation back into development. \S\ref{sec:delphi} puts the proposed
indicators to clinicians before any model is fitted, because no amount of
statistics repairs an indicator that measures the wrong thing.
\S\ref{sec:plausibility} asks whether the three domains correspond to known
mechanisms in haemophilia, since a model can predict well and still be wrong
about why. \S\ref{sec:hemo_systems} reads the haemophilia model of
Appendix~\ref{app:hemo_dynamics} as a control system and asks where it fails,
which is where the general systems view of Chapter~\ref{ch:systems} meets a
disease whose mechanism is written down in equations. \S\ref{sec:hemo_roadmap}
sets out the five phases of external validation that would settle the question
with data, and \S\ref{sec:hemo_close} states what the chapter settles and what
it leaves to them.

Profile~B of Table~\ref{tab:profiles} serves as the reference case. Its
biomedical coordinate is the same in every jurisdiction; its structural
coordinate is not, and the profile moves from $(0.30,0.70,0.85)$ where
treatment is funded to $(0.30,0.70,0.20)$ where it is rationed. The technical
passages below are read against that pair of values.
\end{ataglance}

\begin{quote}
\small\itshape
Note to the reader. This chapter is conceptual and forward-looking; no haemophilia data are analysed here. The figure and derivation blocks below are reserved in place and marked as placeholders, to be filled once the disease-specific study is carried out. This note is removed in the final version.
\end{quote}

\noindent The clinical picture comes first, because it is the one the haemophilia clinic already knows. Haemophilia is an inherited deficiency of clotting factor~VIII or factor~IX. Without treatment it brings repeated bleeding into joints and muscles, and each joint bleed adds to a damage that does not fully reverse: synovial inflammation, loss of cartilage and, in time, a fixed arthropathy that limits movement. Factor replacement given as prophylaxis, longer-acting concentrates, emicizumab and now gene therapy have changed what the disease means \citep{srivastava_2020}, so that many people with haemophilia now grow old with it and meet the hypertension, diabetes, heart and kidney disease of age on top of their joint damage \citep{mauser_2009}. The decisions that shape their course are familiar ones: which product and dose to give, how often to measure factor activity and inhibitor titre, whether to start immune tolerance induction when an inhibitor appears, whether to offer gene therapy. Each of these decisions assumes that product arrives, that the person can reach a treatment centre and that someone at home can help give treatment. Two people with the same factor level and the same joint score can therefore have very different futures. The rest of this chapter asks whether the model built in the earlier chapters can describe that difference, and the model enters only after this picture.

Can a measure of health capacity built in one population be carried into a disease it never saw and still mean the same thing? The chapters before this one answered narrower questions. Chapter~\ref{ch:results1} asked whether the three domains of the $H_{\text{potential}}$ can be measured in a general cohort of older adults, and Chapter~\ref{ch:results2} asked whether the resulting score predicts death and first overnight hospital stay in that cohort. Both questions were put to the population in which the construct was built. A construct that works only where it was built describes that cohort; it has not yet been shown to describe health. This chapter asks whether the construct survives being moved, and takes haemophilia as the first place to move it. It addresses Aim~4.

The question matters now for two reasons. The first is general. Prediction models moved to a new population often perform worse there, because differences in patients, health systems and the frequency of the outcome change the link between predictors and outcomes \citep{moons_2019,collins_2024}. If the $H_{\text{potential}}$ is to be used anywhere outside the cohort that produced it, the thesis must say how it would be moved and how the move would be checked, before anyone tries. The second reason is specific to haemophilia. Treatment has turned the disease, for many people, from a fatal childhood condition into a lifelong one \citep{srivastava_2020}, and the factors that decide long-term outcome have widened with that change. A measure that weighs biology, social resources and the health system together has a use in such a disease that it would not have had when survival was the only question.

The $H_{\text{potential}}$ was built in a general ageing cohort precisely so that it would not depend on any one disease before it was adapted to one. Translation from that cohort to haemophilia is therefore not a matter of swapping one list of clinical variables for another. It requires a judgement about whether the latent construct still means the same thing once it is seen through the biology, treatment and social setting of a particular condition. The previous chapters developed and tested the general framework; this chapter asks how the framework can be specialised without changing what it is meant to measure. Its aim is not to build a new model for haemophilia. It is to show how a construct that does not depend on any one disease can be carried into a disease-specific setting while keeping the assumptions that motivated it.

The difference between development and translation shapes the whole thesis. Many disease-specific outcome measures are built inside one clinical population. They may work well there, but what they measure becomes tied to the signs and treatment of that disease, so they are hard to compare across conditions. The $H_{\text{potential}}$ was built the other way round: in a mixed ageing population first, with any disease-specific use deferred. The aim of that order is that the construct should describe a general property of health and not the natural history of one disorder.

Building the framework in a cohort defined by age rather than by disease has three effects, and each can be checked. It makes it less likely that the latent structure leans on features unique to one disease, since no single disease dominates the cohort. It lets the domains of Biomedical Potential, Social Modifiers and Systemic Burdens form around determinants of health shared across chronic conditions. And it leaves a fixed measurement structure that can later be specialised by replacing indicators while the construct stays the same. Translation in this sense does not alter the theoretical model. It changes the observed variables used to estimate the model in a given clinical setting. This follows the general rule of translational medicine, that the mechanism must be kept while the way it is measured is adapted to the target population.

The alternative would have been to build a haemophilia measure from the start, in a haemophilia cohort. That route was rejected because it would have answered a different question. A measure built in haemophilia would be likely to centre on factor level and joint damage, and it would then have no way to say whether the social and systemic domains carry information of their own or merely stand in for the severity of bleeding. Building first in a general cohort, where those domains were defined without reference to bleeding, keeps that question open and testable.

This chapter is therefore a conceptual and methodological bridge between the general framework of the previous chapters and disease-specific empirical work still to come. No haemophilia cohort is analysed in this thesis, and no disease-specific measurement model is estimated. The chapter instead sets out the case for translation, names clinically meaningful indicators for each domain, describes how expert consultation would establish their content and face validity, and proposes a staged plan for external validation and recalibration in an independent haemophilia cohort. What it offers are commitments made before the data exist, so that the data, when they come, can show those commitments to be wrong.

The approach is deliberately cautious. Following current guidance on external validation and recalibration, it does not assume that the factor loadings, thresholds or prognostic relationships of the ageing cohort carry over to haemophilia. It keeps the latent structure and re-estimates the parameters in the target population.

Three things could go wrong in a translation of this kind, and none of them would announce itself. The name of a domain could be kept while its meaning shifts, so that Biomedical Potential in haemophilia comes to mean ``how severe the bleeding disorder is'' and nothing more. Scores from the two cohorts could be set side by side as though they lay on one scale when they lie on two. And a threshold set in the ageing cohort could be carried over without re-estimation, so that everyone in the new cohort is misclassified in the same direction. The sections below guard against each of them in turn, and the chapter says plainly where a guard depends on data that do not yet exist.

The chapter proceeds in the order set out at its head. Section~\ref{sec:why_haemophilia} gives the reasons for choosing haemophilia. Section~\ref{sec:hemo_mapping} maps the domains onto haemophilia indicators and states the cost of the substitution, and Section~\ref{sec:two_stage} shows that the two-stage logic survives it. Section~\ref{sec:delphi} puts the indicators to expert consultation, and Section~\ref{sec:plausibility} asks whether each domain matches a known mechanism. Section~\ref{sec:hemo_systems} reads the haemophilia model as a control system and uses it to test the general systems view of Chapter~\ref{ch:systems}. Section~\ref{sec:hemo_roadmap} sets out the five phases of external validation, and Section~\ref{sec:hemo_close} closes the chapter.

% =====================================================================
\section{Why Haemophilia}
\label{sec:why_haemophilia}
% =====================================================================

\noindent The opening of the chapter left one choice unexplained: why the first disease should be haemophilia and not another. A first translation has to be informative whichever way it turns out. A disease in which outcome is set almost wholly by biology would make the social and systemic domains look idle whatever their true worth; a disease whose biology is poorly understood would leave no way to tell whether the biomedical domain had been measured at all. Haemophilia sits between these. Its central biological defect is known and measurable, and yet outcome in it depends heavily on whether treatment reaches the person. That combination makes it a test the framework can fail.

Haemophilia also shows clearly how chronic disease has changed in the present century. It was long marked by repeated spontaneous bleeding, progressive damage to muscles and joints, and much shorter life. Before effective clotting factor replacement, severe bleeding was often fatal, and care aimed mainly at stopping acute bleeds and prolonging survival. Treatment has since changed the picture. Prophylactic factor replacement, factor concentrates with a longer half-life, non-factor therapies such as emicizumab, and the arrival of gene therapy have turned haemophilia, for many people, from a fatal childhood disease into a lifelong condition \citep{srivastava_2020}.

As survival has improved, the aims of care have changed with it. Care now goes well beyond preventing fatal bleeding. It covers keeping muscles and joints healthy, keeping people able to live on their own, quality of life, mental health, taking part in school and work, and managing the chronic diseases of age that now affect older people with haemophilia \citep{mauser_2009}. Long-term health is therefore no longer decided by the clotting defect alone. It follows from how the severity of that defect combines with access to treatment, the health services available, social support, health literacy, money and the wider health system. This is the same shift, from present disease to future capacity for health, that motivates the $H_{\text{potential}}$.

The shift matches the starting point of the framework. Throughout this thesis, health has been treated not as the absence of disease but as the capacity to keep physical, functional and social wellbeing in the face of biological and environmental challenge. Haemophilia tests that idea well because outcome depends more and more on how medical treatment interacts with the systems in which people live. People whose disease looks equally severe on biological grounds may follow very different courses over the years, depending on their access to specialist care, whether they keep to prophylaxis, their economic resources and whether supply of treatment holds. Factor activity alone cannot explain that spread. If it could, the framework would have nothing to add, and the test would show it.

The literature on haemophilia care increasingly takes this wider view. Current models of care rest on a team of haematologists, orthopaedic surgeons, physiotherapists, nurses, psychologists, social workers and rehabilitation staff, and accept that ageing well with haemophilia means managing the biological, functional, psychological and social sides of health together. The three domains of the $H_{\text{potential}}$ make the same claim.

The point is sharpest in low- and middle-income countries, where large gaps remain in access to diagnosis, prophylaxis and team-based care. Even where effective therapies exist, breaks in the supply of medicines, distance from specialist centres, lack of money and unequal health services may change long-term outcome a great deal. The same biological severity may therefore lead to very different clinical histories depending on where care is given. By placing Social Modifiers and Systemic Burdens beside Biomedical Potential, the framework gives a way to record these wider determinants within one model. In a setting where the supply of treatment is uncertain, that is not a refinement; it is a source of variation the measure must capture.

Haemophilia also tests how far the $H_{\text{potential}}$ generalises. The construct was built in a mixed ageing population whose health varied across many body systems and chronic conditions. A successful translation would show that the construct is not confined to the setting in which it was built, and that it captures a wider property of health which can be measured through different disease-specific indicators. Haemophilia is complex enough to press the framework hard. Long-term outcome in it depends on inherited severity, joint damage built up over years, keeping to treatment, the development of inhibitors, ageing, other chronic disease, psychological wellbeing, access to care and the arrival of new therapies. A framework that can hold these interacting causes together without losing its meaning would have a strong claim to use in other chronic conditions where biology, health systems and social setting are similarly entangled.

One further reason is particular to this thesis. Appendix~\ref{app:hemo_dynamics} has already written the core of haemophilia treatment as a small set of equations: factor activity rising with infusion and falling with clearance and with inhibitor neutralisation, the inhibitor titre growing and being suppressed by immune tolerance induction, and gene therapy adding a source of factor that the inhibitor also neutralises. Few chronic diseases allow the treatment loop to be written down so plainly. This matters for the systems layer of the thesis. Chapter~\ref{ch:systems} makes general claims about stocks, delays, thresholds and control; haemophilia is a disease in which those claims can be checked against a known mechanism rather than argued in general terms (\S\ref{sec:hemo_systems}). For these reasons haemophilia was chosen as the first translational application of the $H_{\text{potential}}$.

Choosing the disease is the easy half. The harder question is what, in
practice, each domain is measured by once the general indicators are set aside.

% =====================================================================
\section{Mapping the Domains to Haemophilia Indicators}
\label{sec:hemo_mapping}
% =====================================================================

\noindent The previous section gave reasons to try haemophilia; it did not say what would be measured once there. The premise of the translation is that the latent construct stays the same across clinical settings, while the observed indicators used to estimate it change to suit the biology and health care of a particular disease. Translation therefore happens at the level of measurement, not of theory. The three domains found in the general ageing population, Biomedical Potential, Social Modifiers and Systemic Burdens, keep their definitions. What changes is how each domain is represented in data, through haemophilia indicators that have both biological relevance and clinical face validity. The value of the framework in haemophilia is that it gives weight to the social and systemic factors that modify the biological disease, so that outcome is no longer treated as a straight-line function of clotting factor level.

The distinction matters. If translation required redefining the latent construct, each disease would need its own theory, comparison across conditions would be impossible, and the aim of a measure of health capacity that does not belong to one disease would fail. The approach here assumes instead that the determinants of future health act through common mechanisms that show themselves differently in different clinical settings. In psychometric terms, this means holding the construct fixed while allowing its indicators to be substituted. In confirmatory factor analysis a latent variable is not observed; it is inferred from several correlated indicators. Those indicators need not be the same in every population, provided they measure the same underlying construct and meet the usual requirements of reliability and validity. Translation to haemophilia therefore requires indicators that are clinically meaningful, consistent with the definitions of the latent domains, and able to support estimation of the measurement model in a future haemophilia cohort.

The proposed mapping is set out in Table~\ref{tab:hemo_mapping} and drawn in Figure~\ref{fig:hemo_domains}. The table lists, for each domain, the indicators proposed and the direction in which each is expected to move the score; the figure starts from the person ageing with haemophilia, groups the indicators into three boxes, one per domain, and shows each domain feeding the $H_{\text{potential}}$, the first two raising it and Systemic Burdens, drawn dashed, lowering it until it is reverse coded. The right of the figure lists the outcomes the score would be asked to predict in Phase~4 of the roadmap: recurrent bleeding, progression of arthropathy, unplanned hospitalisation, orthopaedic intervention, loss of quality of life and death from any cause. Read from left to right, the figure is the chapter's claim in one line: the clinic's own measures, grouped by domain, should predict the outcomes the clinic cares about. Table~\ref{tab:hemo-map} then sets each indicator used in the ageing cohort beside the haemophilia indicator proposed in its place. The indicators were chosen from a review of the haemophilia literature, current international clinical guidelines, and outcome measures routinely collected in haemophilia care programmes. Three rules governed the choice. An indicator had to be something clinics already record, since an indicator that needs a new data collection would delay validation by years. It had to sit clearly inside one domain, since an indicator that belongs equally to two would blur the boundary the measurement model relies on. And it had to have a stated direction, since an indicator whose effect on future health could run either way cannot be scored. The list is representative, not complete. Its purpose is to show how the domains may be measured in haemophilia, not to fix an instrument. Final selection is deferred to the Delphi consultation described later in the chapter.

\begin{table}[htbp]
\centering
\caption{Mapping the $H_{\text{potential}}$ domains to haemophilia-specific indicators.}
\label{tab:hemo_mapping}
\footnotesize
\renewcommand{\arraystretch}{1.25}
\begin{tabular}{p{2.8cm} p{6.3cm} p{2.8cm}}
\toprule
\textbf{Domain} & \textbf{Representative haemophilia indicators} & \textbf{Expected direction} \\
\midrule
Biomedical Potential & Joint health (Haemophilia Joint Health Score \citep{hilliard_2006}); annualised bleed rate; bleeding phenotype; functional mobility; inhibitor status; age-related comorbidity & Better joint and functional health $\rightarrow$ higher $H_{\text{potential}}$ \\
\addlinespace
Social Modifiers & Access to a haemophilia comprehensive-care centre; distance and travel time to specialist care; caregiver support; peer support; health literacy; treatment self-management & Greater support and access $\rightarrow$ higher $H_{\text{potential}}$ \\
\addlinespace
Systemic Burdens & Treatment supply-chain stability; factor availability; out-of-pocket treatment costs; insurance or funding security; structural barriers to care & Greater burden $\rightarrow$ lower $H_{\text{potential}}$ (reverse coded) \\
\bottomrule
\end{tabular}
\end{table}

\begin{figure}[htbp]
\centering
% Chapter 8 (validation.tex): the three domains for a person ageing with haemophilia.
% Print-safe: domains differ by grey shade and by border style; burden path is dashed.
{\hpdiagramsizes
\begin{tikzpicture}[node distance=3.5mm, font=\footnotesize]
\tikzset{
  dom/.style={draw, rounded corners=2pt, align=left, text width=5.9cm, inner sep=4pt, minimum height=1.55cm},
  bpd/.style={dom, fill=black!4, thick},
  smd/.style={dom, fill=black!14, thick, densely dotted},
  sbd/.style={dom, fill=white, thick, dashed},
}
\node[hpnote, text width=5.9cm, align=left, font=\footnotesize\itshape] (who)
  {A person ageing with haemophilia: inherited deficiency of factor~VIII or IX,
   joint damage built up over years, a possible inhibitor, and the comorbidity of age};
\node[bpd, below=of who] (bp) {\textbf{Biomedical Potential}\\
  Joint health (HJHS); annualised bleed rate; bleeding phenotype; functional
  mobility; inhibitor status; age-related comorbidity};
\node[smd, below=of bp] (sm) {\textbf{Social Modifiers}\\
  Access to a comprehensive-care centre; distance and travel time to specialist
  care; caregiver and peer support; health literacy; treatment self-management};
\node[sbd, below=of sm] (sb) {\textbf{Systemic Burdens}\\
  Supply-chain stability; factor availability; out-of-pocket costs; insurance or
  funding security; structural barriers to care};
\node[hpkey, text width=1.55cm, minimum height=1.6cm, right=0.95cm of sm] (h) {$H_{\text{potential}}$};
\node[hpbox, align=left, text width=3.55cm, font=\footnotesize, right=1.25cm of h] (out)
  {\textbf{Outcomes of interest}\\[1pt]
   Recurrent bleeding\\ Progression of arthropathy\\ Unplanned hospitalisation\\
   Orthopaedic intervention\\ Loss of quality of life\\ All-cause mortality};
\draw[hparrow] (bp.east) -- node[above, pos=0.45] {$+$} (h.north west);
\draw[hparrow, densely dotted] (sm.east) -- node[above] {$+$} (h.west);
\draw[hparrow, dashed] (sb.east) -- node[below, pos=0.45] {$-$} (h.south west);
\draw[hparrow] (h) -- node[above=1pt, font=\scriptsize] {predicts} (out);
\node[hpnote, font=\scriptsize\itshape, text width=13.6cm, below=3mm of sb.south west, anchor=north west]
  {Border and arrow style mark the domain: solid, Biomedical Potential; dotted, Social Modifiers;
   dashed, Systemic Burdens, which lowers the score and is reverse coded before estimation.
   No haemophilia data are analysed; indicators are proposals for the Delphi panel.};
\end{tikzpicture}}

\captionsetup{justification=justified, width=\linewidth}
\caption{The three domains of the $H_{\text{potential}}$ as they apply to a person ageing with haemophilia. Each box lists the example indicators proposed in Table~\ref{tab:hemo_mapping}; the right-hand box lists the outcomes of interest named for prognostic evaluation (\S\ref{sec:hemo_roadmap}). Border and arrow style distinguish the domains, so the figure reads in black and white.}
\label{fig:hemo_domains}
\label{fig:hemo-domains}
\end{figure}

\begin{table}[htbp]
\centering
\caption{Indicators of the ageing cohort and the haemophilia indicators proposed in their place, by domain. General indicators are those of Table~\ref{tab:loadings}. ``Direct'' marks a pair that measures the same property; ``partial'' a pair that overlaps; ``none'' an indicator with no counterpart on the other side.}
\label{tab:hemo-map}
\footnotesize
\renewcommand{\arraystretch}{1.2}
\begin{tabular}{p{2.3cm} p{3.4cm} p{4.6cm} p{1.6cm}}
\toprule
\textbf{Domain} & \textbf{General indicator (NHATS)} & \textbf{Haemophilia counterpart proposed here} & \textbf{Match} \\
\midrule
Biomedical Potential & Multimorbidity & Age-related comorbidity & Direct \\
 & ADL/IADL limitation & Functional mobility & Partial \\
 & Cognitive function & No counterpart proposed & None \\
 & Psychological wellbeing & No counterpart proposed & None \\
 & No counterpart & Joint health (HJHS); annualised bleed rate; bleeding phenotype; inhibitor status & None \\
\addlinespace
Social Modifiers & Support network & Caregiver support; peer support & Direct \\
 & Communication & Health literacy & Partial \\
 & Social participation & No counterpart proposed & None \\
 & No counterpart & Access to a haemophilia care centre; distance and travel time; treatment self-management & None \\
\addlinespace
Systemic Burdens & Economic vulnerability & Out-of-pocket treatment costs; insurance or funding security & Partial \\
 & Healthcare barriers & Structural barriers to care & Direct \\
 & Educational attainment & No counterpart proposed & None \\
 & Environmental stress & No counterpart proposed & None \\
 & No counterpart & Supply-chain stability; factor availability & None \\
\bottomrule
\end{tabular}
\end{table}

Table~\ref{tab:hemo-map} makes plain how little the two indicator sets share. Every haemophilia indicator stays in the domain of the general indicator nearest to it, which is what Condition~2 of Assumption~\ref{ass:hem:substitution} below requires. But only a few pairs measure the same property, and the indicators that carry most of the clinical weight in haemophilia, the joint score, the bleed rate, inhibitor status and the supply of product, have no counterpart in the ageing cohort at all. The pairings are the author's reading of the domain definitions, not a result, and the Delphi panel of \S\ref{sec:delphi} is free to change them. The table is the reason the derivation and proposition that follow conclude that scores from the two cohorts cannot be set on one scale.

The mapping follows the structure set out across the thesis. Biological determinants remain part of future health, but they are no longer the whole of it. Biological capacity interacts with social circumstances and with the health system to decide whether physical reserve becomes good long-term outcome. This suits haemophilia, where the same degree of clotting factor deficiency may lead to very different clinical courses depending on access to treatment, adherence, rehabilitation and health services.

The indicators in Table~\ref{tab:hemo_mapping} are not all of one kind, and the difference matters for how the domain is read. In the language of \S\ref{sec:sys-stocks}, some indicators measure a stock and some a flow. The Haemophilia Joint Health Score records the state of the joints at the time of assessment, which is the sum of every bleed that has damaged them and every repair that has not undone the damage. It is a stock, and it changes slowly. The annualised bleed rate records how many bleeds occur per year, which is a flow into that stock. A person with a poor joint score and a low bleed rate has a large stock of damage and a small current inflow: well treated now, after years of poor treatment. A person with a good joint score and a high bleed rate has the reverse: little damage so far and a fast inflow. The two would contribute similar amounts to a domain score built from a simple sum, yet their futures differ, because the second person's stock is still filling. The measurement model of Chapter~\ref{ch:results1} treats each indicator as a reflection of one latent level and has no term for the difference between level and rate. The point is noted here so that it is not lost: if the stock and flow indicators load weakly together in a future haemophilia cohort, that would be the first explanation to examine, before concluding that the domain is poorly measured.

The next block states, in the notation of the measurement model, what is kept and what is lost when the indicators change. The model it uses is the one a laboratory uses to read an assay. A domain score cannot be observed directly; it is inferred from several indicators, as an analyte concentration is inferred from an assay signal through a calibration. In the equation below, $\mathbf x_g$ is the list of indicator values recorded for a person in cohort $g$, the ageing cohort $N$ or the haemophilia cohort $H$. The vector $\boldsymbol\tau_g$ holds the intercepts, the reading each indicator gives at a domain score of zero, which plays the part of a calibration offset. The matrix $\boldsymbol\Lambda_g$ holds the factor loadings, the slope that says how strongly each indicator tracks its domain, as a calibration slope says how strongly a signal tracks the analyte. The vector $\boldsymbol\xi_g$ holds the three domain scores, and $\boldsymbol\varepsilon_g$ the measurement error of each indicator. The point the derivation makes is that a new set of indicators is a new assay, and a new assay must be calibrated in its own right.

\begin{derivation}[Construct invariance under indicator substitution]
\label{der:indicator_substitution}
For the developmental and haemophilia settings, write
\[
 \mathbf x_g=\boldsymbol\tau_g+\boldsymbol\Lambda_g\boldsymbol\xi_g+
 \boldsymbol\varepsilon_g,
 \qquad g\in\{N,H\}.
\]
Indicator substitution preserves the intended construct only if the same domain
definitions and loading-zero pattern apply (configural \citep{putnick_bornstein_2016} invariance), orientations
are aligned, and the latent variables share an admissible common scale. Equality
of $\boldsymbol\Lambda_N$ and $\boldsymbol\Lambda_H$ is not automatic when
different indicators are used, so conventional item-level metric invariance is
generally untestable without anchor indicators. Translation therefore establishes
content correspondence, not numerical equivalence. Comparable scores require
shared anchors, an explicit linking model, or local re-estimation followed by
external calibration; changing $\mathbf x$ and $\boldsymbol\Lambda$ while merely
renaming $\boldsymbol\xi$ does not prove construct invariance.
\end{derivation}

\begin{reading}[What indicator substitution does and does not buy]
The derivation makes a narrow point with wide consequences. Calling a domain
Biomedical Potential in both cohorts keeps the name; it does not keep the
ruler. The biomedical score of about $0.55$ for Profile~A comes from the
ageing indicators of Table~\ref{tab:loadings}: chronic conditions, daily
function, cognition and psychological wellbeing. The $0.30$ for
Profile~B would come from joint health, bleed rate, inhibitor status and
mobility. No indicator is common to the two, so the two numbers do not lie on
one scale, and the statement that Profile~B has less biomedical potential than
Profile~A has no support in the model.

What the numbers do support is position within a cohort. Profile~B at $0.30$
can be placed against other adults with haemophilia, and Profile~A at $0.55$
against other older adults. Making the two comparable would need indicators measured in both
cohorts, or a linking study designed for the purpose. Neither exists here, and
the chapter does not pretend otherwise.
\end{reading}

The chapter's central claim is that the framework travels. What follows is the
price of that claim, stated as three conditions that can each be checked and
each fail. A laboratory that replaces one assay with another for the same
analyte asks three matching questions: does the new method measure the same
analyte, does its signal rise when the analyte rises, and has it been
calibrated afresh rather than read off the old method's curve? The parallel
breaks at one point. A laboratory can run both methods on the same samples and
compare them directly; here no person is measured on both indicator sets, so
the comparison of methods that a laboratory would make is not available.

\begin{assumption}[What indicator substitution requires]
\label{ass:hem:substitution}
Let the developmental and haemophilia measurement models be
$\mathbf x_g=\boldsymbol\tau_g+\boldsymbol\Lambda_g\boldsymbol\xi_g+
\boldsymbol\varepsilon_g$ for $g\in\{N,H\}$, with
$E(\boldsymbol\varepsilon_g)=\mathbf 0$. Substituting the general indicators
$\mathbf x_N$ with the haemophilia indicators $\mathbf x_H$ of
Table~\ref{tab:hemo_mapping} preserves the construct only if all three of the
following hold.
\begin{enumerate}
  \item \textbf{The domains keep their meaning.} $\boldsymbol\xi_N$ and
  $\boldsymbol\xi_H$ are defined by the same three domain definitions, in the
  same order and with the same orientation, so that a rise in a coordinate
  means the same thing in both settings. Nothing but the observable variables
  is allowed to change.
  \item \textbf{The new indicators load on the same factors.} Each column of
  $\boldsymbol\Lambda_H$ has its non-zero entries in the same domain block as
  the indicator it replaces, and no indicator loads on a domain it was not
  assigned to. Written the other way, the loading-zero pattern of
  $\boldsymbol\Lambda_H$ equals that of $\boldsymbol\Lambda_N$ and every
  $\lambda_{jk}\ne0$ has the sign the domain definition implies. Equality of
  the non-zero values themselves is not assumed and, with disjoint indicator
  sets, is not testable.
  \item \textbf{The scale is re-identified.} The metric of
  $\boldsymbol\xi_H$ is fixed afresh in the haemophilia cohort, by a reference
  loading or by setting the factor variance to one, and the choice is
  reported. A score carries no meaning until this has been done, and it
  carries no cross-cohort meaning even then.
\end{enumerate}
\end{assumption}

\begin{reading}[Which of the three can actually fail]
The first fails silently. A panel that quietly reads Biomedical Potential as
``severity of haemophilia'' has changed the construct while keeping the label,
and no fit statistic will report it. The Delphi consultation of
\S\ref{sec:delphi} exists to catch this one, because it is the only condition
here that statistics cannot test.

The second fails loudly and is the easiest to check: fit the model and look at
where the indicators land. Health literacy loading on Biomedical Potential
would say the assignment is wrong.

The third is not a failure so much as a step that is easy to skip. Skip it and
the haemophilia score has no stated units, and the $0.30$ of Profile~B and
the $0.55$ of Profile~A appear to share one scale when they lie on two.
\end{reading}

The three domains are taken one at a time below, starting with the biological
one, because its haemophilia indicators are the best established and the other
two are easiest to judge against them.

\subsection{Biomedical Potential}

Biomedical Potential is the person's own physical capacity to maintain health and recover from the effects of disease. In the ageing cohort it was measured through physical reserve, functional ability and general biological health. In haemophilia the same ideas apply but must be measured through the signs the disease leaves on the body.

The principal indicator proposed for this domain is joint health, assessed using the Haemophilia Joint Health Score (HJHS). Chronic haemarthroses remain one of the defining pathological consequences of haemophilia, and repeated bleeding episodes lead progressively to synovial inflammation, cartilage destruction, chronic pain, reduced range of motion and permanent arthropathy. The HJHS provides a standardised clinical assessment of joint impairment and is widely regarded as one of the most informative measures of long-term musculoskeletal health in people with haemophilia. Because preservation of joint integrity directly influences mobility, independence and quality of life, joint health represents an appropriate indicator of underlying biological capacity.

Annualised bleed rate (ABR) provides a complementary measure of disease activity. Whereas joint health reflects the cumulative consequences of previous bleeding episodes, the annualised bleed rate captures ongoing haemostatic stability. Lower bleed rates indicate improved disease control and reduced exposure to progressive joint damage. However, ABR should not be interpreted in isolation because bleeding frequency is influenced by treatment adherence, access to prophylaxis and lifestyle factors in addition to intrinsic disease severity. Bleeding phenotype further refines biological characterisation by recognising that individuals with similar factor levels frequently exhibit different clinical bleeding patterns, since genetic modifiers, musculoskeletal status, treatment adherence and behavioural factors all contribute to observed phenotype.

Functional mobility constitutes another essential indicator. Although mobility is influenced partly by environmental and social factors, it remains closely linked to cumulative biological integrity in haemophilia because progressive arthropathy directly impairs physical function. Inhibitor status represents a further clinically important indicator: development of inhibitors against replacement clotting factor substantially complicates treatment, reduces therapeutic effectiveness and increases the risk of recurrent bleeding and progressive joint damage. Finally, age-related comorbidity becomes increasingly important as life expectancy improves, since older adults with haemophilia now experience hypertension, diabetes mellitus, cardiovascular disease, chronic kidney disease and other conditions that interact with haemophilia-specific pathology to influence functional reserve. Taken together, these indicators provide a multidimensional representation of Biomedical Potential that extends beyond laboratory measures of clotting factor activity.

One omission from the list is deliberate. Baseline clotting factor level, the measure by which haemophilia severity is usually classed, is not proposed as an indicator of the domain. The reason lies in the test the framework must eventually pass. Phase~4 of the roadmap (\S\ref{sec:hemo_roadmap}) asks whether the score predicts future bleeds and admissions beyond what factor level and joint score already predict. If factor level were built into the score, that comparison would be partly circular: the score would contain the comparator, and some of any apparent gain would be the comparator predicting itself. Keeping factor level outside the domain leaves it free to serve as the benchmark. The cost is that the domain is measured through the consequences of the defect rather than the defect itself, which is consistent with what Biomedical Potential is meant to be: capacity, not diagnosis.

Two of the indicators carry a risk that the others do not. Annualised bleed rate and inhibitor status both depend on treatment as well as on biology: the bleed rate falls when prophylaxis is given and kept to, and, as \S\ref{sec:plausibility} shows, whether a given inhibitor titre is protected depends on the dose. Both indicators may therefore load partly on Systemic Burdens or Social Modifiers as well as on Biomedical Potential. Condition~2 of Assumption~\ref{ass:hem:substitution} forbids such cross-loadings in the specified model. If they appear in data, the honest response is to report them and to ask whether the indicator should be measured differently, for example by recording bleed rate together with the regimen in force, rather than to free the loading quietly and keep the fit.

None of these indicators records whether the person can act on the capacity
it measures. The same joint damage has different consequences for a person who
lives twenty minutes from a haemophilia centre and for one who lives four
hours from it.

\subsection{Social Modifiers}

The second domain recognises that biological capacity alone does not decide future health. People with the same physical status often follow very different courses because they differ in their ability to use health services, keep to treatment and obtain support. Social Modifiers are the resources and circumstances that let biological capacity become good outcome.

Access to a haemophilia comprehensive care centre is proposed as the principal indicator within this domain. Comprehensive care centres provide integrated multidisciplinary management including haematology, physiotherapy, orthopaedic care, nursing, psychosocial support and patient education, and numerous studies have demonstrated improved clinical outcomes among individuals managed through such programmes compared with fragmented models of care. Geographical accessibility remains particularly important in many healthcare systems, since distance from specialist centres and travel time influence attendance at scheduled appointments, continuity of prophylactic treatment, participation in rehabilitation and timely management of acute bleeding episodes. These factors are especially relevant in South Africa and other low- and middle-income countries where specialised haemophilia services are often concentrated within major urban centres.

Caregiver support and family involvement constitute additional indicators, because lifelong management of haemophilia frequently depends upon support with treatment administration, transport, appointment attendance and emotional wellbeing. Peer support similarly contributes to long-term disease management through educational resources, practical advice and psychosocial support, particularly during major therapeutic transitions such as initiation of prophylaxis or gene therapy \citep{srivastava_2020}. Health literacy represents another key determinant, because understanding the nature of haemophilia, recognising early bleeding symptoms, administering treatment appropriately and navigating complex healthcare systems all require substantial knowledge and self-management capacity. Together these indicators record that future health depends on the social resources available to turn biological capacity into function, and not on that capacity alone.

The domain has a weakness that should be stated before data expose it. Access to a care centre and distance from it are partly properties of where a person lives, and where a person lives also governs much of what Systemic Burdens records, such as whether local supply is stable. The two domains may therefore be strongly correlated in a haemophilia cohort, more so than in the ageing cohort, where social and systemic indicators were drawn from different parts of daily life. A high correlation between factors is not in itself a failure: the framework expects the domains to be related. But if the correlation approaches one, the data cannot tell the two domains apart, and the three-domain structure would collapse to two. The Delphi panel will be asked to judge whether travel time belongs with Social Modifiers, as proposed, or with the structural barriers of the third domain, and the measurement model of Phase~2 will report the factor correlation so that the question can be answered with data.

Social resources are broadly things a person or a family can hold and
build. The next domain covers the constraints they can neither hold nor
negotiate away, which is why it needs separate accounting.

\subsection{Systemic Burdens}

The third domain captures barriers imposed by the healthcare system itself. Unlike Biomedical Potential and Social Modifiers, which describe capacities and resources, Systemic Burdens quantify constraints that reduce the ability of individuals to benefit from available treatments. These factors sit largely outside the patient's control and still shape long-term outcomes heavily.

Treatment supply-chain stability represents perhaps the most distinctive systemic determinant within haemophilia care. Modern management is critically dependent upon uninterrupted access to clotting factor concentrates, non-factor replacement therapies or emerging gene therapies, and interruptions in procurement, manufacturing or distribution may result in treatment interruption despite otherwise optimal clinical management. Closely related is treatment availability, since access to prophylaxis differs substantially across healthcare systems and economic settings. Financial barriers also warrant explicit consideration: out-of-pocket treatment costs, insurance limitations and funding insecurity may influence treatment adherence, attendance at specialist clinics and willingness to seek medical care during bleeding episodes. Broader structural barriers, including fragmented referral pathways, administrative delays, inconsistent reimbursement policies and limited availability of specialised rehabilitation services, further influence the extent to which effective treatment can be delivered.

Because higher values of this domain represent greater disadvantage, all indicators are reverse coded before estimation of the latent construct. Following reverse coding, higher values consistently represent greater health potential across all three domains, simplifying interpretation of the resulting latent score.

Reverse coding is more than a convenience of sign. In the realisation model of Chapter~\ref{chap:crf}, the burden score $B$ enters as structural freedom $R=1-B$, and $R$ multiplies with the social score inside the factor that converts biomedical potential into realised capability. Worked Example~\ref{ex:hem:mrb} below shows what that means in practice: with biology and social support held fixed, a change in $R$ alone cuts realised capability by more than half. If the domain were left with its original orientation, a high score would mean a large burden, and a model that multiplied it in would reward burden. The orientation must therefore be fixed before estimation and checked after it, by confirming that each burden indicator loads with the sign the reverse coding implies. An indicator whose sign comes out wrong is more likely to have been coded wrongly than to have revealed a new fact about haemophilia.

The domain also differs from the other two in timing. A supply interruption changes a person's Systemic Burdens score as soon as it happens, while the joint damage it causes builds over years. Systemic Burdens is the fastest-moving of the three domains and Biomedical Potential the slowest. A single measurement captures the burden at one moment and the damage accumulated from all earlier moments, so the two are measured on different time scales even when they are recorded on the same day. \S\ref{sec:hemo_systems} returns to this difference, because it is where delays enter.

Reverse coding puts the three domains on a common orientation. It says
nothing about how they act on one another, and that is where the case for a
single latent construct actually rests.

\subsection{Integration of the Three Domains}

Although presented separately for conceptual clarity, the three domains are unlikely to operate independently, and their principal strength lies in their interaction. Excellent Biomedical Potential cannot be fully realised if Systemic Burdens prevent access to treatment. Similarly, optimal healthcare systems may achieve only limited benefit when severe biological impairment has already occurred. Social Modifiers mediate many of these relationships by determining whether patients are able to engage effectively with available healthcare resources and sustain appropriate long-term self-management. This interaction represents one of the principal conceptual advances of the $H_{\text{potential}}$: rather than treating biological disease, social context and healthcare systems as competing explanations for health outcomes, the framework views them as complementary determinants of a single latent capacity. The proposed mapping should consequently be regarded as a theoretically grounded starting point for empirical model development rather than a final measurement instrument.

The worked example below puts numbers on this interaction for one reference case, Profile~B of Table~\ref{tab:profiles}: an adult with haemophilia, a high-titre inhibitor and damaged joints, scored once where treatment is funded and once where it is rationed. It uses the realisation model of Chapter~\ref{chap:crf}. In words, that model takes biomedical potential $P$ and scales it by a factor between zero and one set jointly by social support $S$ and structural freedom $R$; with the weight $\alpha=0.5$ used here, the factor is the square root of $S$ times $R$. Realised capability $H$ is then the share of biomedical potential that the person's circumstances allow to be used. For the clinician, the example asks a single question: how much of the same patient's capacity is lost when only the funding of treatment changes?

\begin{workedexample}[hem:mrb]{Realised capability of Profile~B from haemophilia indicators}
Score Profile~B on the indicators of Table~\ref{tab:hemo_mapping}. A poor
Haemophilia Joint Health Score, a high annualised bleed rate, a high-titre
inhibitor and limited mobility put Biomedical Potential at $P=0.30$. Access to
a comprehensive-care centre, caregiver support with treatment and good health
literacy give Social Modifiers $S=0.70$. Systemic Burdens are scored from supply-chain stability, funding
security and out-of-pocket cost, then reverse coded to structural freedom
$R=1-B$. Under a funded regimen $R=0.85$; where bypassing agents are rationed
and specialist care is distant, $R=0.20$. The scores are
illustrative, since no haemophilia cohort is analysed here.

With $\alpha=0.5$ the funded multiplier is $\sqrt{0.70\times0.85}=0.771$ and
the rationed one is $\sqrt{0.70\times0.20}=0.374$. Realised capability is
$H=0.30\times0.771=0.231$ funded and $H=0.30\times0.374=0.112$ rationed. As a
fraction of biomedical potential, $H$ is $77$ per cent under the funded regimen
and $37$ per cent under the rationed one, and the ratio between the two states is $2.06$.

\begin{wecode}[Python]
from hpv.crf import H, contextual_mean

P, S = 0.30, 0.70
for label, R in (("funded", 0.85), ("rationed", 0.20)):
    mult = contextual_mean(S, R, rho=0.0)
    print(f"{label:9s} multiplier {mult:.3f}  H {H(P, S, R):.3f}")

print(H(P, S, 0.85) / H(P, S, 0.20))    # 2.06
\end{wecode}

\weanswer{$H=0.231$ funded and $H=0.112$ rationed, that is $77$ and $37$ per
cent of biomedical potential, a ratio of $2.06$. Asserted by
\texttt{test\_worked\_example\_hem\_profile\_b\_favourable\_and\_unfavourable}.}
\end{workedexample}

The example is small, but it carries the argument of the chapter. Biomedical potential and social support are the same in both states; only structural freedom differs. An additive index would move by the change in $R$ times its weight and no more, and the size of that move would not depend on $P$ or $S$. In the multiplicative form the loss from rationing is a share of what the person could otherwise realise, so the same fall in $R$ costs more in absolute terms when biology and social support are strong. A clinic that improved Profile~B's joint health without any change in funding would therefore find its gain scaled down by the rationed multiplier of $0.374$, not by the funded one of $0.771$. This is the precise sense in which the framework says the domains interact: the return on improving any one of them depends on the level of the others. It is also a claim that data could refute, since a cohort in which the effects of the domains on outcome simply add would favour the additive index over the realisation model.

A mapping of this kind is defensible only if the sequence that produced it is
intact: general construct first, disease-specific indicators second. What that
sequence commits the work to, and what it rules out, is set out next.

% =====================================================================
\section{Preserving the Two-Stage Logic}
\label{sec:two_stage}
% =====================================================================

\noindent A defining characteristic of the proposed framework is its separation of construct development from disease-specific application. This two-stage logic underpins the methodological integrity of the thesis and distinguishes the present work from many disease-specific measurement studies. Rather than developing a latent construct directly within a single clinical population, the framework was first conceptualised, operationalised and evaluated in a heterogeneous population of older adults before any attempt was made to adapt it to an individual disease. Translation therefore represents a second stage in the research programme rather than a continuation of model development. The overall design is shown in Figure~\ref{fig:two_stage}.

\begin{figure}[htbp]
\centering
\begin{tikzpicture}[node distance=10mm]
\node[hpbox, text width=5.6cm, minimum height=2.6cm] (s1) {\textbf{Stage 1: general ageing cohort}\\[2pt] Develop and validate the disease-independent construct in NHATS\\[2pt]\footnotesize Chapters~\ref{ch:methods}--\ref{ch:results2}};
\node[hpbox, text width=5.6cm, minimum height=2.6cm, right=of s1] (s2) {\textbf{Stage 2: haemophilia cohort}\\[2pt] Substitute disease-specific indicators, re-estimate and recalibrate\\[2pt]\footnotesize Chapter~\ref{ch:translational}};
\draw[hparrow] (s1) -- (s2);
\node[hpkey, minimum height=0.9cm, text width=12.4cm, below=7mm of $(s1.south)!0.5!(s2.south)$, anchor=north]
  (fixed) {Latent domains held fixed: Biomedical Potential, Social Modifiers, Systemic Burdens};
\draw[hpdash] (s1.south) -- (s1.south |- fixed.north);
\draw[hpdash] (s2.south) -- (s2.south |- fixed.north);
\end{tikzpicture}

\captionsetup{justification=justified, width=\linewidth}
\caption{Two-stage design separating construct development from disease-specific translation.}
\label{fig:two_stage}
\end{figure}

This sequence is intentional. Measurement models derived exclusively from a single disease population inevitably risk incorporating characteristics that are unique to that disease rather than representing broader determinants of health. While such models may exhibit strong performance within their original clinical context, their applicability beyond that setting is often uncertain. By contrast, the $H_{\text{potential}}$ was designed from the outset to measure a general latent property of health that is independent of any particular diagnosis. The ageing cohort therefore functions as the developmental population in which the theoretical structure of the construct is established, while disease-specific cohorts serve as independent populations in which that structure is tested, refined and, where necessary, recalibrated. This mirrors contemporary recommendations for prediction-model research and psychometric instrument development, which consistently emphasise that model development, internal validation and external validation represent distinct phases of scientific investigation rather than interchangeable analytical procedures \citep{moons_2019,collins_2024,steyerberg_2019}.

The haemophilia adaptation therefore reuses the domain structure and the two-stage design of Chapters~\ref{ch:methods} to~\ref{ch:results2}, replacing the general indicators with the disease-specific indicators of Table~\ref{tab:hemo_mapping} and re-estimating the loadings and thresholds in a haemophilia cohort. If translation required altering the latent domains themselves, the resulting framework would represent a different construct rather than a disease-specific application of the original model. Instead, only the observable indicators are expected to change. Biomedical Potential remains Biomedical Potential, Social Modifiers remain Social Modifiers, and Systemic Burdens remain Systemic Burdens. What differs is the manner in which these latent concepts are measured. From a structural equation modelling perspective, this corresponds to maintaining the latent factor structure while permitting disease-specific indicator substitution: measures such as daily function or chronic disease burden used within the ageing cohort are replaced by haemophilia-relevant variables such as joint health, annualised bleed rate and treatment accessibility, while the hypothesised relationships among the latent domains remain unchanged.

Figure~\ref{fig:two_stage} shows this as two stages joined by a fixed layer. Stage~1 is the general ageing cohort of Chapters~\ref{ch:methods} to~\ref{ch:results2}; Stage~2 is the haemophilia cohort of this chapter; and the layer between them is the set of three latent domains, which the figure marks as held fixed. Read in the terms of \S\ref{sec:sys-boundary}, the latent domains form an interface. On one side sit the indicators, which may be replaced; on the other sit the realisation model and the prognostic analysis, which consume domain scores and do not need to know which indicators produced them. An engineered system that is built this way can have a component replaced without the rest being redesigned, provided the new component meets the same interface specification. Assumption~\ref{ass:hem:substitution} is that specification. Its three conditions say what a set of haemophilia indicators must satisfy to be plugged in behind the domains without breaking anything downstream.

The interface view also makes clear what would go wrong if the sequence were reversed or merged. Suppose the domains were redefined during translation, for instance by moving inhibitor status into Systemic Burdens because inhibitor management depends on product supply. The realisation model downstream would still run and still return numbers, but $P$ and $R$ would no longer mean what they meant when the model was specified, and every interpretation that rested on those meanings, such as the reading of Worked Example~\ref{ex:hem:mrb}, would fail without any error being raised. A change to an interface that is not declared as such is the kind of fault that passes every local check. The two-stage design guards against it by forbidding domain redefinition in Stage~2 outright: if the haemophilia data show that the domains do not hold, the correct conclusion is that translation has failed, not that the domains should be quietly adjusted until it succeeds.

Maintaining this separation between latent structure and observable indicators also facilitates comparison across diseases. If subsequent adaptations are developed for other chronic conditions, such as sickle cell disease, Parkinson's disease or chronic kidney disease, each may employ its own set of clinically relevant indicators while retaining the common latent framework established in this thesis. Another advantage is that the strategy permits recalibration without conceptual modification. It is anticipated that factor loadings, indicator thresholds and latent score distributions will differ between the general ageing population and a haemophilia cohort because the prevalence, distribution and clinical meaning of individual indicators are unlikely to be identical. Such differences do not imply that the theoretical model has failed; rather, they reflect expected variation arising from differences in disease biology, healthcare systems and patient characteristics. Re-estimation of measurement parameters is therefore viewed as recalibration rather than redevelopment of the framework, mirroring developments in contemporary prognostic modelling where externally validated models are frequently recalibrated for new populations without altering their underlying conceptual structure.

Measurement invariance is the property an assay has when it gives the same reading for the same concentration in different patient groups or in different laboratories. The derivation below lists the four standard levels of that property, each stricter than the last, and says which comparison each level permits. It applies where the same indicators are recorded in both groups, which is the favourable case; the reading after Proposition~\ref{prop:hem:scalar} restates the levels in plain words and says how few of them can be tested here.

\begin{derivation}[Measurement invariance across populations]
\label{der:invariance}
With common indicators in groups $g=1,2$, configural invariance imposes the same
loading-zero pattern; metric invariance adds
$\boldsymbol\Lambda_1=\boldsymbol\Lambda_2$; scalar invariance adds
$\boldsymbol\tau_1=\boldsymbol\tau_2$; and strict invariance additionally equates
residual variances. Nested models are compared using estimator-appropriate scaled
$\chi^2$ differences together with changes in practical fit indices and inspection
of local misfit. Metric invariance supports comparison of relationships involving
the factors; scalar invariance is normally required for latent-mean or absolute
score comparisons. If only partial invariance holds, at least two invariant anchors
per factor and transparent freeing of non-invariant parameters are required.
\end{derivation}

The level that decides whether two cohort means may be set side by side is
scalar invariance, and the reason is a counting argument rather than a
statistical one. The laboratory form of the argument is familiar. Two
laboratories report different mean values for the same analyte in their
patients. The gap may reflect a true difference between the patient groups, a
difference in calibration offset between the two methods, or both, and from the
reported values alone the two causes cannot be told apart. They can be told
apart only by running a common reference material in both laboratories. The
proposition shows that cohort means of the $H_{\text{potential}}$ behave in the
same way: the intercepts $\boldsymbol\tau$ play the part of the calibration
offset, the latent means $\boldsymbol\kappa$ the part of the true concentration,
and indicators shared by both cohorts, called anchors, the part of the
reference material.

\begin{proposition}[A raw difference of cohort means is uninterpretable without scalar invariance]
\label{prop:hem:scalar}
For $g\in\{N,H\}$ let $\mathbf x_g=\boldsymbol\tau_g+\boldsymbol\Lambda
\boldsymbol\xi_g+\boldsymbol\varepsilon_g$ with $E(\boldsymbol\varepsilon_g)=
\mathbf 0$, a shared $p\times q$ loading matrix $\boldsymbol\Lambda$ of full
column rank, and latent means $\boldsymbol\kappa_g=E(\boldsymbol\xi_g)$. The
observed means satisfy
\[
  E(\mathbf x_H)-E(\mathbf x_N)
  =\underbrace{(\boldsymbol\tau_H-\boldsymbol\tau_N)}_{\boldsymbol\delta}
   +\boldsymbol\Lambda(\boldsymbol\kappa_H-\boldsymbol\kappa_N).
\]
If $\boldsymbol\delta$ is left free, then for every candidate difference
$\Delta\boldsymbol\kappa$ there is a $\boldsymbol\delta$ that reproduces the
observed means exactly, so $\Delta\boldsymbol\kappa$ is not identified. Scalar
invariance, $\boldsymbol\delta=\mathbf 0$, identifies it and gives
$\Delta\boldsymbol\kappa=(\boldsymbol\Lambda^{\mathsf T}\boldsymbol\Lambda)^{-1}
\boldsymbol\Lambda^{\mathsf T}\{E(\mathbf x_H)-E(\mathbf x_N)\}$. Absent that
restriction, or an equivalent set of anchor indicators, the difference between
a haemophilia cohort mean $H_{\text{potential}}$ and an NHATS cohort mean
carries no interpretation as a difference in capability. The statement assumes
a shared $\boldsymbol\Lambda$, which is the favourable case. Where the two
indicator sets are disjoint, as in Table~\ref{tab:hemo_mapping}, the loadings
are not comparable either, and the conclusion holds all the more strongly.
\end{proposition}

The proof counts what the data can supply against what has to be found. Each
indicator supplies one observed mean difference; each intercept difference and
each latent mean difference is an unknown. When the unknowns outnumber the
equations, more than one answer fits the data equally well.

\begin{proof}[Proof sketch]
The mean structure supplies $p$ equations. The unknowns are the $p$ intercept
differences and the $q$ latent mean differences, so the system is
under-determined by $q$ and the mean structure is saturated. Fix any
$\Delta\boldsymbol\kappa$ and set $\boldsymbol\delta=E(\mathbf x_H)-
E(\mathbf x_N)-\boldsymbol\Lambda\Delta\boldsymbol\kappa$: the fit is unchanged
and the latent means have moved. A shift of intercepts is therefore
observationally indistinguishable from a shift of latent means. Setting
$\boldsymbol\delta=\mathbf 0$, or fixing at least $q$ of its entries through
anchors common to both cohorts, removes the $q$ excess degrees of freedom and
leaves the stated least-squares solution \citep{meredith_1993}.
\end{proof}

\begin{reading}[The four levels of invariance, in plain words]
Each level asks a harder question of the same data. Configural: do the same
indicators group into the same three domains in both cohorts? Metric: does each
indicator move by the same amount when the domain moves? Scalar: does a person
with the same true capability give the same expected answer in both cohorts?
Strict: is the leftover noise the same size?

Only the first is likely to be testable here, since metric invariance needs
indicators shared between the cohorts and the haemophilia set shares almost
none of them. The consequence is a limit on what may be said. Suppose the mean
biomedical score is $0.34$ in the haemophilia cohort and $0.51$ in the ageing
cohort. Without scalar invariance that gap of $0.17$ cannot be read: it may be
a real difference in capability, or the joint health score may simply start its
scale in a different place from the ageing cohort's activity count. Relationships survive where
levels do not, so a finding that a one-standard-deviation fall in the domain
raises admission risk by a given amount can be set beside the same finding in
the other cohort long before the two means can.
\end{reading}

The proposition and its reading set a firm limit on the claims the chapter can make. Any sentence of the form ``people with haemophilia have lower health potential than older adults in general'' is ruled out, however natural it may sound, because no set of indicators shared by both cohorts exists to fix the common scale. What can be said is relational: whether the same domain predicts the same kind of outcome with a similar strength, and whether the domains relate to one another in the same pattern. These are weaker statements, but the design supports them, and stating the limit now keeps the stronger claim from entering by accident later. A study that wanted the stronger claim would need to measure a few indicators in both kinds of cohort, such as a common measure of mobility, and use them as anchors in the sense of Derivation~\ref{der:invariance}. That is a separate study and is not proposed here.

Finally, preserving the two-stage logic enhances scientific transparency. Separating development from translation allows each stage of the research to be evaluated independently. The validity of the latent construct does not depend upon the success or failure of any single disease application, and the quality of a disease-specific adaptation can be assessed without revisiting the development of the original framework. Such modularity facilitates cumulative scientific progress because future adaptations can build upon an established measurement architecture while remaining open to empirical refinement.

Modularity is a claim about process, and a process can be honoured on paper
while the indicator set itself is wrong. Before anything is estimated, the
indicators have to be put to people who treat haemophilia.

% =====================================================================
\section{Establishing Face and Content Validity: A Modified Delphi Consultation}
\label{sec:delphi}
% =====================================================================

\noindent Before any disease-specific measurement model can be estimated, it is necessary to establish that the proposed indicators are clinically meaningful and appropriately represent the latent domains they are intended to measure. Although the indicator mapping presented in the previous section is grounded in published literature and existing haemophilia outcome measures, theoretical justification alone is insufficient to establish content validity. Consequently, a structured programme of expert consultation is proposed prior to empirical estimation of the translated measurement model. Face validity refers to the extent to which a proposed measure appears, on expert judgement, to assess the construct it claims to measure; content validity extends this concept by evaluating whether the complete set of indicators adequately represents the conceptual domain of interest. While these forms of validity are inherently qualitative, they constitute an essential foundation for subsequent psychometric evaluation because statistical analyses cannot compensate for poorly chosen or conceptually inappropriate indicators.

To establish face and content validity, this thesis proposes the use of a modified Delphi consultation involving clinicians and allied health professionals associated with the Wits Haemophilia Comprehensive Care Centre. The Delphi method is widely employed in health research when empirical evidence is incomplete or when expert judgement is required to refine complex conceptual frameworks, and its principal strength lies in its ability to develop structured consensus while minimising the influence of dominant individuals through anonymous and iterative feedback. A modified Delphi design was selected rather than a classical Delphi because the starting point is not an empty list of candidate indicators. Instead, the indicator set has already been developed through theoretical reasoning and a review of the haemophilia literature, and the role of the expert panel is therefore to evaluate, refine and prioritise these proposed indicators rather than to generate an entirely new conceptual framework. The process is shown in Figure~\ref{fig:delphi}.

\begin{figure}[htbp]
\centering
\begin{tikzpicture}[node distance=5mm]
\tikzset{st/.style={hpbox, text width=6.6cm, minimum height=0.95cm}}
\node[st] (p) {Proposed indicator set};
\node[st, below=of p] (r1) {Round 1: experts rate each indicator for relevance, clarity and direction};
\node[st, below=of r1] (fb) {Anonymised, controlled feedback of the group ratings};
\node[st, below=of fb] (rn) {Further rounds of rating};
\node[draw, diamond, aspect=2.6, fill=black!6, font=\small, align=center, inner sep=1pt, below=of rn] (c) {Pre-specified\\consensus reached?};
\node[hpkey, text width=6.6cm, minimum height=0.95cm, below=6mm of c] (out) {Indicators retained, revised or dropped};
\draw[hparrow] (p) -- (r1);
\draw[hparrow] (r1) -- (fb);
\draw[hparrow] (fb) -- (rn);
\draw[hparrow] (rn) -- (c);
\draw[hparrow] (c) -- node[right, font=\small] {yes} (out);
\draw[hparrow, dashed] (c.east) -- ++(1.3,0) |- node[pos=0.25, right, font=\small] {no} (fb.east);
\end{tikzpicture}

\captionsetup{justification=justified, width=\linewidth}
\caption{Modified Delphi consultation used to establish face and content validity of the haemophilia indicator set.}
\label{fig:delphi}
\end{figure}

The proposed panel would comprise experienced clinicians and allied health professionals involved in the long-term management of people with haemophilia. Representation from multiple professional disciplines is important because the three latent domains extend beyond biomedical management alone: haematologists provide expertise regarding disease severity and treatment strategies; physiotherapists contribute knowledge relating to musculoskeletal function and rehabilitation; specialist nurses provide insight into treatment adherence and patient education; and psychologists and social workers contribute perspectives on psychosocial wellbeing and social support. The multidisciplinary composition of the panel therefore mirrors the multidimensional structure of the $H_{\text{potential}}$ itself.

The consultation is expected to proceed through two or three iterative rounds. During the first round, panellists would independently evaluate each proposed indicator with respect to conceptual relevance, clinical importance, clarity of definition and appropriateness of assignment to its proposed latent domain, and would indicate the expected direction of association between each indicator and the $H_{\text{potential}}$. Following completion of the first round, responses would be summarised anonymously and returned to participants together with aggregated descriptive statistics. Controlled feedback is one of the defining characteristics of the Delphi methodology because it enables experts to reconsider their previous judgements in light of the collective opinion of the panel without introducing direct interpersonal pressure. Subsequent rounds would focus on achieving convergence of opinion rather than introducing substantial conceptual change, and the exact number of rounds is guided by the stability of panel responses rather than by a predetermined schedule. Consensus thresholds should be specified before commencement of the consultation to minimise subjective interpretation of panel responses; although no universal standard exists, thresholds of approximately 70--80\% agreement are commonly employed in health research.

The consultation is defined by the following elements.

\begin{itemize}
  \item \textbf{Panel.} Clinicians and allied health professionals with experience in the long-term care of people with haemophilia.
  \item \textbf{Rounds.} Two to three rounds in which panellists rate each proposed indicator for relevance and clarity on a Likert scale, and comment on the expected direction of each indicator on the $H_{\text{potential}}$.
  \item \textbf{Consensus.} A pre-specified agreement threshold determines which indicators are retained, revised, or dropped, with controlled feedback between rounds.
  \item \textbf{Output.} A refined, clinically endorsed indicator set for each domain, with agreed directions, ready for empirical estimation in a haemophilia cohort.
\end{itemize}

The rule below turns ``consensus'' into numbers that are fixed before any rating is seen, in the way a laboratory fixes its quality-control acceptance limits before a run rather than after it. For each proposed indicator it uses two quantities: $A_j$, the share of panellists who rate the indicator 3 or 4 on a four-point relevance scale, and $\operatorname{IQR}_j$, the interquartile range of the ratings, which measures how far the panel is spread. High agreement with a narrow spread retains an indicator; low agreement removes it; the band between sends it back for revision.

\begin{derivation}[Delphi consensus and stability criteria]
\label{der:consensus}
For item $j$ rated on a four-point relevance scale, define
\[
 A_j=\frac{1}{K}\sum_{k=1}^{K}\mathbb I(r_{jk}\geq3),
 \qquad \operatorname{IQR}_j=Q_{0.75}(r_j)-Q_{0.25}(r_j).
\]
Before round one, retain an item if $A_j\geq0.80$ and
$\operatorname{IQR}_j\leq1$, revise it if $0.70\leq A_j<0.80$, and remove it if
$A_j<0.70$, subject to documented review of essential minority concerns. Between
rounds, define stability as $|A_j^{(r)}-A_j^{(r-1)}|<0.10$ with no category change.
Stop after two consecutive stable rounds or the pre-specified maximum of three.
These thresholds must be declared before ratings are examined; consensus supports
content validity but does not establish psychometric validity.
\end{derivation}

\begin{reading}[The consensus rule on a panel of twelve]
Take a panel of twelve: four haematologists, two physiotherapists, two
specialist nurses, two social workers and two psychologists. Each rates every
proposed indicator from 1 to 4 for relevance.

Travel time to the specialist centre is rated 3 or 4 by ten of the twelve, so
$A_j=10/12=0.83$, and the middle half of the ratings spans one point, so
$\operatorname{IQR}_j=1$. Both conditions hold and the indicator is retained.
Peer support is rated 3 or 4 by eight, giving $A_j=0.67$, below the removal
threshold. It is dropped unless a panellist has recorded a concern the rest of
the panel had not considered, in which case the concern is written down and the
item goes to a further round.

Stability is a separate test. If travel time scores $0.83$ in round one and
$0.88$ in round two, the change is $0.05$ and both values sit in the retain
band, so the item has settled. A move from $0.83$ to $0.78$ is the same size
but crosses into the revise band, so it has not. Fixing these numbers before
any rating is seen is the point of the exercise; chosen afterwards, they would
simply read the panel's answers back to it.
\end{reading}

Three features of the design deserve a defence, because each could have been chosen otherwise.

The first is anonymity between rounds. An open meeting of the same twelve people would be faster and would allow argument. It was rejected because the panel is unequal by design. In the panel of twelve used above, four haematologists sit with two members of each other discipline, and in an open meeting the larger and more senior group would be expected to set the terms. The domains most at risk of being underweighted are the social and systemic ones, which are the domains the smaller disciplines know best. Anonymous rating with controlled feedback lets a social worker's judgement on caregiver support count as much as a haematologist's, which is the reason the panel is mixed at all.

The second is the source of the panel. All panellists are drawn from one centre, the Wits Haemophilia Comprehensive Care Centre. This makes the consultation feasible, but it has a cost that should be named. Clinicians who work together share a local practice, a local patient population and local habits of thought, and their agreement is therefore not twelve independent judgements. In the language of \S\ref{sec:sys-fmea}, a single-centre panel is open to common-cause failure: a blind spot shared by the centre will appear as consensus, and the consensus rule will retain it. The remedy is not available within this thesis, but its form is clear. A second panel at a centre with a different funding environment, rating the same indicators against the same pre-declared thresholds, would show which judgements travel and which belong to one centre.

The third is the absence, in the design as proposed, of people living with haemophilia. The panel consists of clinicians and allied health professionals, and the indicators it rates include several, such as caregiver support, health literacy and travel time, on which patients and families hold knowledge that clinicians see only in part. Whether a parallel patient panel should be added, and how its ratings should be combined with the clinical panel's, is left open here. It is listed among the questions to settle before the consultation begins, since the consensus rule must be declared before any rating is seen.

There is also a limit that no design of the panel removes. The Delphi consultation can establish that experts agree an indicator is relevant and correctly assigned. It cannot establish that the indicator behaves in data as the experts expect. A panel may agree that health literacy belongs in Social Modifiers, and the fitted model may still find it loading weakly there because it is poorly recorded. Agreement and measurement are different properties, tested at different stages, and the second can overturn the first.

The principal output of the Delphi process is not a validated measurement instrument but a refined set of clinically endorsed candidate indicators judged to possess satisfactory face and content validity. This refined indicator set provides the foundation for the subsequent psychometric stage of the research programme, during which confirmatory factor analysis, reliability assessment and measurement invariance \citep{meredith_1993,vandenberg_lance_2000} testing would be undertaken in an independent haemophilia cohort. The Delphi consultation therefore occupies a distinct position within the overall research programme: it neither replaces empirical validation nor constitutes evidence that the translated framework possesses satisfactory measurement properties, but is an intermediate step linking theoretical development to quantitative evaluation.

Expert agreement establishes that the indicators are worth measuring. It does
not establish that the mechanisms behind them behave in haemophilia the way the
framework assumes, and that has to be argued from what is known of the disease.

% =====================================================================
\section{Biological Plausibility}
\label{sec:plausibility}
% =====================================================================

\noindent A central requirement of any translational framework is that the proposed construct possesses biological plausibility within the disease to which it is being applied. Biological plausibility refers to the extent to which the relationships implied by a conceptual model are consistent with established knowledge of disease mechanisms, clinical progression and therapeutic response. Although statistical association is necessary for developing useful prediction models, prediction alone is insufficient if the underlying conceptual relationships contradict accepted biological understanding. The biological rationale for the present framework differs from traditional disease models in one important respect: conventional prognostic models often begin with the disease itself and attempt to identify variables associated with subsequent outcomes, whereas the $H_{\text{potential}}$ begins with the individual's capacity to maintain health, so that disease becomes one determinant of that capacity rather than its sole defining feature. Within haemophilia, this perspective is increasingly supported by contemporary clinical evidence, since current outcomes reflect successful long-term disease management, preservation of musculoskeletal function, adherence to prophylaxis, psychological wellbeing and management of age-related comorbidity rather than catastrophic haemorrhage alone \citep{srivastava_2020}.

The argument is made one domain at a time, then for the interaction among
them, which is where the multiplicative form earns or loses its place.

\subsection{Biomedical plausibility}

The first domain, Biomedical Potential, represents the individual's intrinsic physiological reserve and capacity to maintain health despite chronic disease. In haemophilia this reserve is influenced by multiple interacting mechanisms, including baseline disease severity, cumulative musculoskeletal damage, recurrent bleeding, inhibitor development, ageing and multimorbidity. Repeated haemarthroses remain the principal mechanism through which untreated or inadequately treated haemophilia produces long-term disability: recurrent bleeding initiates chronic synovial inflammation, progressive cartilage loss, subchondral bone remodelling and irreversible haemophilic arthropathy, which in turn reduces mobility, physical activity, muscle strength and functional independence, creating a cycle in which reduced movement further accelerates musculoskeletal deterioration. Preserved joint health therefore provides a biologically plausible indicator of retained physiological reserve. Annualised bleed rate reflects the effectiveness of haemostatic control and ongoing exposure to joint injury, while inhibitor development alters the effectiveness of replacement therapy and increases the probability of recurrent bleeding despite otherwise appropriate management. The ageing process strengthens this interpretation further, since older adults with haemophilia increasingly experience cardiovascular disease, diabetes, hypertension, obesity, chronic kidney disease and osteoporosis, which interact with haemophilia rather than replacing it.

Inhibitor status is the one indicator in this domain that can be turned into a
decision the same afternoon, and doing so shows what a Biomedical Potential
score is standing in for.

The clinical question is the one a haematologist faces when a patient on
replacement therapy develops an inhibitor: at what titre, in Bethesda units,
does the current regimen stop keeping factor activity above the level that
prevents bleeding, and how far does a larger dose move that titre? The worked
example answers it with the haemophilia model of
Appendix~\ref{app:hemo_dynamics}. In that model a regimen holds factor activity
at a steady level where infusion balances clearance and neutralisation by the
inhibitor. The quantity $h_{\text{need}}$ below is the largest strength of
neutralisation, as a fraction of its maximum, that the regimen can absorb
before that steady level falls under the protective activity $S_{\text{prot}}$;
the tolerable titre $Y_{\text{thr}}$ is the titre at which neutralisation
reaches that strength. The Hill exponent
$\beta$ sets how sharply neutralisation rises once the titre passes the
half-neutralisation titre $K_I$; with $\beta$ above one, the inhibitor has
little effect at low titre and then a steep one, which matches the clinical
experience that an inhibitor presents as a threshold event rather than a slow
drift (Appendix~\ref{app:hemo_dynamics}). The parameters are illustrative and
in arbitrary units, so the example shows the direction and shape of the effect,
not its clinical size.

\begin{workedexample}[hem:dose]{Moving the inhibitor threshold by changing the dose}
Appendix~\ref{app:hemo_dynamics} fixes a regimen for Profile~B: delivered input
$f+h=60$, neutralisation $\alpha=4$, half-neutralisation titre $K_I=5$~BU,
Hill exponent $\beta=2$ and clearance $\lambda=1$, with a protective activity
of $S_{\text{prot}}=20$. At that input the tolerable titre is $5.0$~BU. The
clinical question is not what the threshold is, but which way it moves when the
dose moves, because the dose is the part the clinic controls.

Raise the delivered input by half, from $60$ to $90$. Then
$h_{\text{need}}=(90/20-1)/4=0.875$ and
$Y_{\text{thr}}=5(0.875/0.125)^{1/2}=13.2$~BU: the tolerable titre has nearly
tripled. Cut the input by a quarter instead, to $45$, and
$h_{\text{need}}=0.3125$ gives $3.4$~BU. The threshold is convex in the dose,
so the last increment buys more than the first: $60$ to $75$ adds $2.4$~BU,
while $75$ to $90$ adds $5.8$.

Two ends of the range have no threshold at all. At an input of $100$ or above
the regimen holds factor activity at or over $S_{\text{prot}}$ whatever the titre, and at
$20$ or below it never reaches $S_{\text{prot}}$ even at zero titre. The code
raises an error in both cases rather than returning a number.

\begin{wecode}[Python]
from hpv.haemophilia import inhibitor_threshold

REG = dict(alpha=4.0, K_I=5.0, beta=2.0, lam=1.0)
for dose in (45.0, 60.0, 75.0, 90.0):
    thr = inhibitor_threshold(numerator=dose, S_prot=20.0, **REG)
    print(f"input {dose:5.1f} -> tolerable titre {thr:5.2f} BU")

try:
    inhibitor_threshold(numerator=110.0, S_prot=20.0, **REG)
except ValueError as e:
    print("no finite threshold:", e)
\end{wecode}

\weanswer{Tolerable titres of $3.37$, $5.00$, $7.42$ and $13.23$~BU as the
delivered input rises from $45$ to $90$. A half increase in dose nearly triples
the tolerable titre. Asserted by
\texttt{test\_worked\_example\_hem\_threshold\_moves\_with\_dose}.}
\end{workedexample}

\begin{reading}[Why this belongs in a measurement chapter]
On the $60$ regimen a titre of $4$~BU is protected and a titre of $8$~BU is
not, yet both are recorded as ``inhibitor positive'' by the indicator in
Table~\ref{tab:hemo_mapping}. The indicator is a summary of a threshold that
depends on the regimen as much as on the person. In this sense Biomedical
Potential in haemophilia is never purely biomedical, because the dose is a
funding decision.

That is an argument for scoring the domain from the titre and the regimen
together where both are recorded, and for saying plainly, where they are not,
that the indicator is coarse in a direction that will attenuate the domain's
association with anything.
\end{reading}

Every mechanism above assumes the treatment reached the person. Whether it
does is not a biological question, and the second domain is where it is asked.

\subsection{Social plausibility}

The second domain, Social Modifiers, reflects the observation that favourable biology alone does not guarantee favourable outcomes. Two individuals with similar disease severity may experience markedly different trajectories because their opportunities to engage effectively with healthcare differ. This proposition is well established within haemophilia care, where successful long-term management depends heavily upon treatment adherence, patient education, caregiver support, participation in comprehensive care programmes and timely recognition of bleeding episodes. Prophylactic replacement therapy, for example, provides maximal benefit only if administered consistently over prolonged periods, and adherence depends on family support, treatment literacy, psychological wellbeing, financial stability and access to care as much as on patient motivation. The importance of multidisciplinary comprehensive care further illustrates the plausibility of this domain, because integrated management involving haematologists, physiotherapists, orthopaedic specialists, specialist nurses, psychologists and social workers is effective precisely because biological treatment alone is insufficient to optimise long-term health.

Support, literacy and adherence describe what happens around the patient.
None of them conjures factor concentrate that the health system has not bought,
which is the separate business of the third domain.

\subsection{Systemic plausibility}

The third domain recognises that health systems affect outcome whatever the characteristics of the patient. Haemophilia shows this clearly because its management depends on continuous access to specialised therapies. Treatment interrupted by medicine shortages, procurement delays, failures of distribution or lack of funding may change outcome substantially while disease severity stays the same, and distance from specialist centres may delay treatment of acute bleeds or reduce preventive care. The effect is clearest across countries: nations with similar patient populations often show very different long-term outcomes because treatment availability, reimbursement, national procurement and specialist workforce differ. In low- and middle-income countries, breaks in the supply of clotting factor, uneven spread of specialist services and the cost of sustained prophylaxis may affect outcome as strongly as biological severity itself.

Each domain is plausible on its own terms. The framework claims more than
that. It claims the domains multiply, and that claim needs its own defence.

\subsection{Interaction between the domains}

Perhaps the strongest argument for the biological plausibility of the proposed framework lies not within the individual domains themselves but within their interaction. Contemporary understanding of chronic disease increasingly rejects linear causal pathways in favour of interconnected systems in which biological, behavioural and environmental processes continuously influence one another. Successful prophylaxis, for example, preserves joint integrity, allowing continued physical activity that maintains musculoskeletal function and reduces future bleeding risk; yet sustained prophylaxis depends upon treatment availability, patient education, caregiver support and stable healthcare systems, so that deficiencies in any one domain may propagate throughout the entire system and ultimately influence long-term health outcomes. The three-domain structure of the $H_{\text{potential}}$ is consistent with this systems perspective and provides the biological rationale for modelling health potential as a latent construct rather than as a simple additive index of clinical variables.

The loops in this account can be named, using the distinction of \S\ref{sec:sys-loops}. One is reinforcing and runs downwards: a bleed damages a joint, the damaged joint limits movement and muscle strength, and reduced movement further speeds the deterioration of muscles and joints, which raises the risk of the next bleed. Left alone, a loop of this kind accelerates, which is why the biomedical plausibility argument above described the damage as a cycle. The other is the same loop run upwards under prophylaxis: protected joints allow continued activity, activity maintains musculoskeletal function, and the risk of future bleeding falls. The two are one structure with opposite signs of drift, and what decides the direction is the input that sits outside the biological domain altogether: whether prophylaxis is supplied and taken. That is the content of the claim that the domains interact. Social Modifiers and Systemic Burdens do not add a separate quantity to biological health; they decide which way the biological loop turns.

This also explains why the plausibility of each domain taken alone is not enough. Each subsection above showed that one domain corresponds to a known mechanism. None showed that the mechanisms combine as the realisation model says they do. A cohort in which the biological loop turned the same way regardless of supply and support would satisfy every domain-level argument and still refute the interaction claim.

The interaction argument has so far been made in words: deficiencies
propagate, loops form, one domain gates another. Appendix~\ref{app:hemo_dynamics}
allows more than that, because it writes the treatment of haemophilia as
equations. The next section reads those equations as a control system, to see
which parts of the verbal argument survive being made exact.

% =====================================================================
\section{A Systems View of Haemophilia Care}
\label{sec:hemo_systems}
% =====================================================================

\noindent The previous section argued in words that the direction of the biological loop is set from outside biology. The general systems view of Chapter~\ref{ch:systems} is also mostly qualitative, because for health in general the mechanism is not known in closed form. Haemophilia is an exception. Appendix~\ref{app:hemo_dynamics} writes the core of its treatment as equations, which can be read as a control system: a process with states, inputs that someone chooses, measurements that someone reads, and feedback through which the choices are revised. This section makes that reading and asks where the loop fails, and whether this one disease bears out the general view or presses against it.

The control vocabulary describes work the haemophilia clinic already does. The process under control is the person's haemostasis. The states are what the clinic would like to know at every moment: factor activity, inhibitor titre and the damage carried in the joints. The inputs are the clinic's decisions: the product and dose of replacement therapy, whether to start immune tolerance induction and how intensely to run it, and whether to give gene therapy. The measurements are what the clinic can actually see: factor activity and titre on the days blood is drawn, bleeds when the person reports them, and the joint score at review. Feedback is the routine of reading those results and changing the dose or the plan. Nothing in this section asks the clinic to work differently. It writes down the loop the clinic already runs, so that its weak points can be named, and so that each weak point can be traced to an indicator of Table~\ref{tab:hemo_mapping}.

\subsection{States, inputs and the process being controlled}

The process has three states that matter here. Circulating factor activity $S_T$ is the first, the quantity a factor assay reports. Read clinically, the result below says that on a fixed regimen and a fixed titre, factor activity does not keep rising or falling but settles at a level set by the dose, the clearance and the inhibitor. It rises with infusion and falls through first-order clearance at rate $\lambda$ and through neutralisation by the inhibitor (equation~\eqref{eq:C_ode}). It is a fast state: for a fixed titre, Proposition~\ref{prop:C_stability} shows that it settles at $S_T^\ast(Y)=(f+h)/(\alpha\,\mathrm{Hill}(Y)+\lambda)$, and that the settling point is stable because clearance grows with the level. This is a balancing loop inside the process itself, and it is why a regimen produces a steady level rather than a rising or falling one.

Inhibitor titre $Y$ is the second state. It has its own dynamics (equation~\eqref{eq:C_iti}): growth that slows as the titre nears a ceiling, natural decay, and suppression by immune tolerance induction. In the model the two states are coupled both ways. The titre lowers factor activity through neutralisation, and factor activity at or above a critical level switches on titre growth, through the indicator term in equation~\eqref{eq:C_iti}. The input that raises $S_T$ therefore also sustains the process that works against it.

Accumulated joint damage is the third state, and in the terms of \S\ref{sec:sys-stocks} it is the stock that matters most. In the stochastic layer of Appendix~\ref{app:hemo_dynamics} it grows whenever factor activity is below the protective level $S_{\text{prot}}$ and does not recover fully. It changes slowly, it integrates every earlier excursion below protection, and it is what the Haemophilia Joint Health Score of Table~\ref{tab:hemo_mapping} measures. Factor activity is the flow-like state that decides how fast the stock fills; joint damage is the stock that Biomedical Potential reads.

The inputs are the levers someone can set. The infused dose $T$ enters through a saturating term, so that doubling a large dose does not double its effect. Immune tolerance induction enters as a suppression rate $I_{\text{ITI}}$ on the titre. Gene therapy enters as a transgene source $A$ that decays with time. These differ in a way that matters for control. The dose can be changed at the next visit. Immune tolerance induction must be held above a critical level for its whole course. Gene therapy is set once, at the moment of infusion, and cannot be adjusted afterwards: in the language of control it is an open-loop input, chosen on a forecast and not corrected by feedback.

\subsection{Measurement, feedback and delay}

A controller acts on what it can measure. In this process the measurements arrive at different speeds and carry different amounts of warning. Titre is measured at clinic visits, so the controller sees $Y$ in samples, not continuously. A bleed is observed when it happens, which is after the protective level has already been crossed. Joint damage is observed through a joint score, which records the sum of past excursions and so arrives last of all. Ordered by the warning they give, titre comes first, bleeds second and joint score last.

The delays matter because of the thresholds. Proposition~\ref{prop:C_threshold} shows that, with a Hill exponent above one, steady-state activity crosses the protective level at a single titre $Y_{\text{thr}}$ and that the crossing is sharp. Worked Example~\ref{ex:hem:dose} puts that titre at $5.0$~BU for the reference regimen. Below it the regimen protects; above it, no change in adherence helps. A process with a sharp threshold, observed through a sampled and lagging measurement, gives the controller little margin. The distance between the current titre and $Y_{\text{thr}}$ is the reserve, in the sense of \S\ref{sec:sys-resilience}: it is the amount of adverse change the system can absorb before its behaviour changes in kind. Two features of that reserve follow from the equations. It belongs to the person and the regimen together, since the threshold moves with the dose. And it can be bought: raising the delivered input from $60$ to $90$ moves the threshold from $5.0$ to $13.2$~BU. The early-warning signal is therefore not a bleed, which marks a crossing that has already happened, but a narrowing gap between titre and threshold, which the model could compute once its parameters had been estimated for a person and regimen.

Immune tolerance induction carries a second threshold, and here delay and threshold combine badly. In clinical terms the question is whether a tolerance course given at a certain intensity will clear the inhibitor or leave it in place. In the model the answer depends on the intensity of suppression $I_{\text{ITI}}$ compared with a critical value set by how fast the titre grows ($\rho$), how fast it decays on its own ($\lambda_Y$) and how strongly the course suppresses it ($\kappa$). By Proposition~\ref{prop:C_iti}, suppression at or above $(\rho-\lambda_Y)/\kappa$ drives the titre to zero, and suppression below it leaves a titre that persists; Worked Example~\ref{ex:apc-iti} shows ninety per cent of the critical suppression leaving a standing titre of $3.6$~BU. A gap in supply part of the way through a course does not cost part of the benefit; it can move the outcome from extinction to persistence. Because the course is long, the controller may learn of the failure well after the supply gap that caused it.

\subsection{Where the loop fails}

The control reading makes it possible to list the failure modes that matter, in the manner of \S\ref{sec:sys-fmea}, and to ask what each would do to the measured score.

\emph{Loss of the actuator.} The dose, the tolerance course and the specialist visit all depend on product and services arriving. A break in supply removes the controller's main input. Because of the thresholds above, the effect is not a small drift in outcome but, possibly, a change of regime: a protected titre becomes unprotected, a tolerance course that would have succeeded leaves a standing inhibitor. In the score it appears first in Systemic Burdens and only later, through joint damage, in Biomedical Potential.

\emph{Saturation of the actuator.} Even with supply intact, the dose has limits. Infusion saturates, and at high titre neutralisation is close to its maximum, so further dose buys little; Appendix~\ref{app:hemo_dynamics} notes that the delivered input is bounded by what can be infused, tolerated and paid for, while the titre is not. The law of requisite variety \citep{ashby_1956} states the general point: a controller can hold a system steady only if it has as many distinct responses as the disturbances it faces. A clinic whose only lever is the dose cannot hold activity against a titre that the dose itself helps sustain. Bypassing agents, which act by a different route, and tolerance induction, which acts on the titre directly, add variety to the controller.

\emph{Coarse sensing.} The inhibitor indicator in Table~\ref{tab:hemo_mapping} records only whether an inhibitor is present, so a protected titre of $4$~BU and an unprotected one of $8$~BU are recorded alike. The stochastic layer of Appendix~\ref{app:hemo_dynamics} adds a second form of coarse sensing. A measure that averages factor activity may class a person as adequately treated while a fraction of each week is spent below protection, and it is that fraction, not the average, that fills the stock of joint damage.

\emph{A wrong model in the controller.} Leveson's account of safety in control structures treats a common cause of loss as a controller whose internal model of the process no longer matches the process \citep{leveson_2011}. The linear dosing arithmetic of Definition~\ref{def:C_linear} is such an internal model. It is the arithmetic used whenever a recovery and a half-life are measured and a dose is calculated from them (Appendix~\ref{app:hemo_dynamics}). Proposition~\ref{prop:C_linear} shows when it is exact; outside that limit it overstates neutralisation, so a dose computed from it errs towards excess. Here the model error falls on the safe side, but that is a property of these equations, not a guarantee, and must be checked again if the model changes.

\emph{Common-cause failure.} The failures above are not independent. A single loss of funding can cut the dose, interrupt a tolerance course and lengthen the wait for specialist care at the same time. Layered defences protect against failure only when the layers do not share a cause \citep{reason_1990}, and in haemophilia care many layers share one: the funding and supply of treatment. For the measurement model this has a direct consequence. When one cause moves several indicators together across domains, the domains will appear more correlated than their definitions imply, and the same cause can produce the cross-loadings that Condition~2 of Assumption~\ref{ass:hem:substitution} forbids. A haemophilia cohort that mixes funded and rationed settings is the place where this is most likely to be seen.

\subsection{The control structure and where the score acts}

The process has more than one controller, and they work on different time scales, as \S\ref{sec:sys-control} describes for health care in general. The innermost loop is the person and caregiver, who give or take prophylaxis, notice bleeds and decide when to seek care; their capacity to run that loop is what Social Modifiers records. Around it sits the treating centre, which sets the dose, reads the titre, and decides on tolerance induction and on whether and when to offer gene therapy. The outermost loop is the funder and the procurement system, which decide what product is bought and whether supply holds.

The feedback each controller receives is unequal. The person sees bleeds as they happen. The centre sees titres and joint scores at clinic visits. The funder sees spending and stock levels, aggregated across patients, and does not see the titre of any one of them. The controller whose decisions set everyone's reserve has the least direct view of it, and this is the gap the $H_{\text{potential}}$ could help to fill. Because the realisation model carries structural freedom $R$ inside the score, a fall in supply appears in the score of each person it touches, before the joint damage it will cause has accumulated. Worked Example~\ref{ex:hem:mrb} shows the size of the effect for Profile~B: realised capability falls from $0.231$ to $0.112$ when $R$ falls from $0.85$ to $0.20$, with biology and social support unchanged. Whether a funder would act on such a signal is a question of governance, not of measurement, and this thesis does not answer it.

\subsection{What the case tests in the general view}

Haemophilia lets several claims of Chapter~\ref{ch:systems} be checked against equations rather than argued. Three of them hold. Capacity behaves as a stock: joint damage fills from the flow of bleeds and does not drain fully, which is the stock-and-flow reading of \S\ref{sec:sys-stocks} in a concrete form. Thresholds are real and depend on the system around the person: both the protective titre and the critical suppression are set partly by the regimen, and so partly by funding. And reserve can be defined exactly, as the gap between present titre and $Y_{\text{thr}}$, which gives the early-warning idea of \S\ref{sec:sys-resilience} a quantity to watch.

One claim is pressed harder than the general view anticipates. The realisation model treats $P$, $S$ and $R$ as separate coordinates measured at one time and combined by multiplication. At a single moment that holds in haemophilia. Over time it does not, because structural freedom reaches into biology: a period of low $R$ lets the titre outrun the dose and lets joint damage accumulate, and that damage is part of the next measurement of $P$. Today's systemic burden becomes a future biomedical deficit. A static score cannot represent that transfer; it can only record its results at each measurement. This is a limit of any static measure applied to a process with memory, not of the haemophilia mapping, and it is the reason the dynamical extension of Appendix~\ref{app:hemo_dynamics} exists as a separate line of work.

The reading has limits of its own, which \S\ref{sec:sys-limits} states in general. The equations of Appendix~\ref{app:hemo_dynamics} use illustrative parameters in arbitrary units and are not fitted to patient data. They show where failures occur and which way a change pushes the process, not how large any effect would be in a real cohort. Every model simplifies the system it stands for \citep{box_1976}; this one at least makes its simplifications explicit enough to be wrong. What the control reading gives the validation programme is a set of things to look for, and answering them needs a cohort, a fitted model and a stated order of work.

% =====================================================================
\section{Roadmap for External Validation and Recalibration}
\label{sec:hemo_roadmap}
% =====================================================================

\noindent The systems reading said what to look for; it did not say in what order, or what result at each step would end the programme. The translational framework proposed in this chapter represents the beginning, rather than the completion, of disease-specific model development. Although the indicator mapping and Delphi consultation establish conceptual and content validity, empirical validation within an independent haemophilia population remains essential before the framework can be considered suitable for clinical application. Contemporary guidance for prediction model research consistently emphasises that externally validating a model in new populations is a critical stage of development rather than an optional refinement \citep{moons_2019,collins_2024}. Accordingly, the present work proposes a structured programme of external validation consisting of five sequential phases, summarised in Figure~\ref{fig:roadmap}.

For the clinical reader, the roadmap is the sequence a new laboratory test goes through before it is used on patients, applied to a score. The indicators are first agreed by experts, as a test's purpose and specimen are agreed before any validation run. The score is then measured in a new population and its structure checked, as the precision and linearity of an assay are checked on local material. It is recalibrated, as a method is calibrated against local reference material. It is then judged on whether it adds information to what the clinic already measures, as a new test is judged on whether it improves diagnosis beyond the tests already in use. Finally the work is approved and reported so that others can repeat it. Each stage produces something that the next one needs, and each ends in a decision that can stop the programme.

\begin{figure}[htbp]
\centering
% Chapter 8 (validation.tex): staged translation roadmap with outputs and decision points.
% Print-safe: phase boxes shaded, outputs plain, decisions bold-bordered; the Delphi stage is dashed.
{\hpdiagramsizes
\begin{tikzpicture}[font=\footnotesize]
\tikzset{
  ph/.style={hpbox, font=\footnotesize, text width=2.75cm, minimum height=1.55cm},
  pr/.style={draw, rounded corners=2pt, align=left, text width=4.65cm, inner sep=3pt, minimum height=1.55cm},
  dc/.style={draw, very thick, align=left, text width=4.3cm, inner sep=3pt, minimum height=1.55cm},
  phd/.style={ph, dashed, fill=white},
  hd/.style={font=\footnotesize\bfseries, align=center},
  link/.style={-{Stealth[length=1.6mm]}, semithick},
}
\def\rowgap{1.95cm}
\foreach \n/\sty/\name/\prod/\dec in {%
  0/phd/{\textbf{Before Phase 1}\\Delphi consultation}/{Clinically endorsed indicator set for each domain, with agreed directions}/{Retain if $A_j\geq0.80$ and $\operatorname{IQR}_j\leq1$; revise if $0.70\leq A_j<0.80$; remove if $A_j<0.70$. Stop after two stable rounds or three},
  1/ph/{\textbf{Phase 1}\\Cohort assembly}/{Registry cohort from comprehensive-care centres, starting in South Africa; route to care recorded}/{Report the observed range of each domain, so a null result can be judged against the cohort boundary},
  2/ph/{\textbf{Phase 2}\\Re-estimation of the measurement model}/{Confirmatory factor model with haemophilia indicators; reliability; factor correlations; configural invariance}/{Cross-loading: revise the indicator and return to the panel. Domains indistinguishable, or wrong signs that survive checking: stop},
  3/ph/{\textbf{Phase 3}\\Recalibration}/{Re-estimated loadings, intercepts and thresholds; risk intercept $a$ and slope $b$; internal validation}/{Set interpretive thresholds anew from the haemophilia distribution; carry none over from NHATS},
  4/ph/{\textbf{Phase 4}\\Prognostic evaluation}/{Association, discrimination, calibration and net benefit against factor level, joint score and inhibitor status}/{No added value: a verdict against clinical use. Added value: check that it comes from the social and systemic domains},
  5/ph/{\textbf{Phase 5}\\Governance, ethics and reporting}/{Ethics approval and data-use agreements, in place before Phase 1; report to TRIPOD}/{Report against the standards of the development study so others can check the answer}%
} {
  \node[\sty] (p\n) at (0,-\n*\rowgap) {\name};
  \node[pr, anchor=west] (o\n) at ($(p\n.east)+(0.4,0)$) {\prod};
  \node[dc, anchor=west] (d\n) at ($(o\n.east)+(0.4,0)$) {\dec};
  \draw[link] (p\n.east) -- (o\n.west);
  \draw[link] (o\n.east) -- (d\n.west);
}
\foreach \a/\b in {0/1,1/2,2/3,3/4,4/5} \draw[hparrow] (p\a.south) -- (p\b.north);
\node[hd, above=1.5mm of p0] {Stage};
\node[hd, above=1.5mm of o0] {What it produces};
\node[hd, above=1.5mm of d0] {Decision before the next stage};
\end{tikzpicture}}

\captionsetup{justification=justified, width=\linewidth}
\caption{Staged roadmap for translating the $H_{\text{potential}}$ into haemophilia. Each row is one stage, in the order it would be attempted; the middle column states what the stage produces and the right-hand column the decision taken before the next stage begins. The Delphi consultation (\S\ref{sec:delphi}), drawn dashed, precedes the five phases of external validation and recalibration.}
\label{fig:roadmap}
\label{fig:hemo-roadmap}
\end{figure}

Figure~\ref{fig:hemo-roadmap} sets out each stage with its product and the decision that follows it. Three of the decisions can end the programme or send it back: the consensus rule of Derivation~\ref{der:consensus} after the Delphi rounds, the failure rules of Phase~2, which return a misassigned indicator to the panel and stop the work if the domains do not reappear, and the test of Phase~4, where a score that adds nothing to factor level, joint score and inhibitor status is a verdict against clinical use. The other decisions protect the reading of later results. Phase~1 must report the range of each domain that the cohort covers, and Phase~3 must set its thresholds afresh. Phase~5 appears last in the list, but its ethics approval and data-use agreements must be in place before Phase~1 begins; only its reporting comes at the end.

At a glance, the roadmap comprises the following steps.

\begin{enumerate}
  \item \textbf{Assemble a haemophilia cohort.} Draw on haemophilia registry and comprehensive-care-centre data, beginning with the South African context.
  \item \textbf{Re-estimate the measurement model.} Fit the domain structure with the disease-specific indicators of Table~\ref{tab:hemo_mapping}, and test measurement invariance against the general-cohort structure.
  \item \textbf{Recalibrate.} Re-estimate loadings, weights, and interpretive thresholds for the haemophilia population.
  \item \textbf{Re-evaluate prognostic value.} Test the recalibrated score against haemophilia-relevant outcomes, such as bleeding events, unplanned hospitalisation, joint deterioration, and mortality.
  \item \textbf{Govern and report.} Obtain ethics approval, observe data-use agreements, and report against the same standards used in the development study.
\end{enumerate}

\begin{reading}[The five phases applied to Profile~B]
Phase 1 enrols adults of the kind Profile~B represents and records them with
the same variables as every other registry member. Phase 2 asks whether joint
score, bleed rate, inhibitor status and mobility behave as one domain in the
data, and
whether the travel and funding items behave as two more. Phase 3 fixes what a
score of, say, $0.23$ means in this population, since it need not mean what
$0.23$ meant in NHATS. Phase 4 asks the question of clinical interest: does
the score, taken at baseline, predict the next joint bleed or the next
admission beyond what factor level and joint score already predict? Phase 5
supplies the ethics approval for the study and the reporting that lets others
check the answer.

If phase 4 returns nothing, the framework has not been shown to be wrong about
haemophilia. It has been shown to add nothing to what haemophilia care already
measures, which for a clinical purpose is the same verdict.
\end{reading}

The phases are set out below in the order they would be attempted. None can be
brought forward, because each assumes an answer to the one before it.

\subsection{Phase 1: Cohort assembly}

The first stage involves establishment of an independent haemophilia cohort. The preferred source is a prospective registry derived from comprehensive care centres, beginning within the South African context and expanding to collaborating centres where feasible. The cohort should include adults across the spectrum of haemophilia severity and treatment modalities, ensuring representation of individuals receiving conventional factor replacement, extended half-life therapies, non-factor replacement therapies and emerging gene therapies where available. Standardised collection of biomedical, psychosocial and healthcare-system variables is essential to preserve comparability across participating centres, and outcome definitions should mirror those employed during development wherever possible while incorporating haemophilia-specific endpoints relevant to long-term disease management.

The source of the cohort sets its boundary, in the sense of \S\ref{sec:sys-boundary}, and the boundary here has a known bias. A registry built from care centres contains people whom care has reached. Those with the heaviest Systemic Burdens, who live furthest from a centre or have lost contact with it, are the least likely to be enrolled and the most likely to be lost to follow-up. The cohort will therefore cover a narrower range of the third domain than exists in the population. A restricted range weakens any association, so a weak link between Systemic Burdens and outcome in such a cohort would not show that the domain does not matter; it might show only that the cohort excludes the people for whom it matters most. Two steps reduce the problem without removing it. Enrolment should record how each person reached the centre and how long they had been known to it, and the analysis should report the observed range of each domain beside its estimated effect, so that a reader can judge whether a null result reflects the construct or the boundary.

A cohort is worth only as much as the structure it can support, and the
first thing to ask of it is whether the three domains reappear at all.

\subsection{Phase 2: Re-estimation of the measurement model}

Following data collection, the haemophilia-specific measurement model should be estimated using confirmatory factor analysis. This stage evaluates whether the three-domain architecture identified within the general ageing cohort is reproduced within the disease-specific population. Internal consistency, convergent validity, discriminant validity and global model fit should be examined using the same methodological framework adopted in the development study, with particular attention to indicator performance because variables demonstrating strong theoretical justification may nevertheless exhibit weak psychometric behaviour within the empirical data. Measurement invariance testing should subsequently evaluate whether the latent structure remains conceptually comparable between the developmental ageing cohort and the haemophilia population, with evidence supporting configural and metric invariance strengthening the interpretation that the same latent construct has been successfully translated rather than replaced.

Because the haemophilia indicators share almost none of the developmental ones, the comparison that Phase~2 can make is mainly configural, as the reading after Proposition~\ref{prop:hem:scalar} explained. The phase therefore needs a stated rule for what counts as failure, set before the model is fitted. Three outcomes would count. If an indicator loads on a domain other than the one it was assigned, Condition~2 of Assumption~\ref{ass:hem:substitution} fails for that indicator. If two domains correlate so highly that they cannot be told apart, the three-domain structure fails in this population. And if the domains appear but with loadings of the wrong sign, the orientation of Condition~1 has failed, most likely through a coding error. Each outcome calls for a different response. The first may be repaired by revising the indicator and returning to the Delphi panel. The second and third, if they survive checking, mean that the construct has not been translated, and the programme should stop at this phase rather than proceed to recalibrate a structure that is not there.

A reproduced structure does not carry the developmental numbers with it.
Loadings, intercepts and thresholds have to be estimated afresh.

\subsection{Phase 3: Recalibration}

Even if the latent structure is preserved, empirical recalibration is expected to be necessary. Differences in disease prevalence, treatment practices, indicator distributions and healthcare environments are likely to alter factor loadings, intercepts and latent score distributions. Recalibration therefore involves re-estimating measurement parameters while maintaining the conceptual relationships among the latent domains, and should be viewed as adaptation to a new population rather than redevelopment of the theoretical framework. Recalibrated scoring algorithms should subsequently be evaluated for stability using internal validation procedures analogous to those employed during the developmental phase. In the development study that check was a bootstrap: the whole analysis is repeated on many samples drawn with replacement from the cohort, much as a validation run is repeated on resampled material, and the spread of results shows how far the scoring would shift in another sample of the same population.

Two different operations go by the name of recalibration, and the derivation below keeps them apart. The first re-estimates how the indicators are turned into a score: the intercepts $\boldsymbol\tau^*$, loadings $\boldsymbol\Lambda^*$ and error variances $\boldsymbol\Theta^*$ of the haemophilia model. In laboratory terms it is the recalibration of a method on local material. The second leaves the score alone and corrects how a score is turned into a predicted risk, through an intercept $a$ and a slope $b$. In laboratory terms it is the check of reported values against known standards, here the observed rates of the outcome in the new cohort. A risk model is well calibrated when its predicted rates agree with the observed ones.

\begin{derivation}[Recalibration of the translated score]
\label{der:recalibration}
Let the translated measurement model be
$\mathbf x^*=\boldsymbol\tau^*+\boldsymbol\Lambda^*\boldsymbol\xi^*+
\boldsymbol\varepsilon^*$. Recalibration re-estimates
$(\boldsymbol\tau^*,\boldsymbol\Lambda^*,\boldsymbol\Theta^*)$ under the retained
factor pattern; the latent covariance should be tested, not automatically fixed.
For a transported prognostic linear predictor $\eta$, calibration at horizon
$t_0$ may be updated as $\eta^*=a+b\eta$, where $a$ corrects baseline risk and
$b$ corrects over- or under-dispersion. Interpretive score thresholds are then
derived anew from the target distribution and decision context. Measurement
recalibration and outcome-risk recalibration are distinct operations: the former
changes how capability is measured, while the latter changes how a fixed score
maps to absolute risk.
\end{derivation}

\begin{reading}[Two different operations both called recalibration]
Suppose Profile~B scores $0.23$ in a jurisdiction that funds treatment. Two
questions follow, and the derivation keeps them apart.

The first is what $0.23$ means. In the ageing cohort the lowest fifth of the
distribution might begin at $0.28$, which would place Profile~B below it. Among
adults with haemophilia and inhibitors the lowest fifth might begin at $0.19$,
which would not. The score has not moved. The reference population has, and the
cut-point with it. Measurement recalibration is the re-estimation of loadings,
intercepts and interpretive thresholds against the new distribution.

The second is what $0.23$ predicts. If the model carried over from the ageing
cohort puts the five-year risk of hospitalisation at 12 per cent and the
haemophilia cohort observes 20 per cent at that score, the intercept $a$
absorbs the difference. If the model spreads high and low scorers too far
apart, the slope $b$ falls below one and pulls predictions back towards the
average. Fixing the risks does not fix the measurement, and fixing the
measurement does not fix the risks.
\end{reading}

The first of those two operations can be put in numbers, and the arithmetic is
short enough that a clinician can check whether a threshold survives the move.
The laboratory parallel is a method with a constant bias. If a new method reads
higher than the old one by a fixed amount, the same reported value now
corresponds to a lower true concentration, and a clinical cut-off that was
fixed on the true concentration has to move on the reported scale. The example
below does the same for one indicator of Biomedical Potential.

\begin{workedexample}[hem:recal]{An intercept shift moves a cut-point}
Take one indicator of Biomedical Potential, standardised so that
$x=\tau+\lambda\xi+\varepsilon$ with $\lambda=0.80$. Suppose Phase~3
re-estimates its intercept in the haemophilia cohort at $\delta=0.072$ above
the developmental value, the loading holding at $0.80$. Every other parameter
is unchanged.

An intercept that is $\delta$ higher means the same observed indicator value is
now produced by a latent score that is $\delta/\lambda$ lower. Any cut-point
anchored to a fixed observed clinical criterion therefore falls by
$0.072/0.80=0.09$. A cut-point set at $\xi=0.28$ in the ageing cohort becomes
$0.19$ in the haemophilia cohort.

Profile~B scores $0.23$. The score lies below the developmental cut-point and
is flagged; it lies above the recalibrated one and is not. The indicators have
not changed and neither has the capability. What changed is the intercept,
which is why Derivation~\ref{der:recalibration} treats the re-estimation of
$\boldsymbol\tau^*$ as part of translation and not as a technical afterthought.

\begin{wecode}[Python]
lam, delta = 0.80, 0.072
cut_dev = 0.28
cut_hem = cut_dev - delta / lam
print(f"cut-point {cut_dev:.2f} -> {cut_hem:.2f}  (shift {delta / lam:.2f})")

profile_b = 0.23
for label, cut in (("developmental", cut_dev), ("recalibrated", cut_hem)):
    print(f"{label:14s} flagged: {profile_b < cut}")
\end{wecode}

\weanswer{The cut-point falls from $0.28$ to $0.19$, and Profile~B at $0.23$ crosses
it in the direction that removes the flag. Asserted by
\texttt{test\_worked\_example\_hem\_intercept\_shift\_moves\_cut\_point}.}
\end{workedexample}

\begin{reading}[What the shift is and is not evidence of]
Nothing here says the haemophilia cohort is healthier. The move of $0.09$ is
the size of a scale artefact, not of a difference in capability, and by
Proposition~\ref{prop:hem:scalar} the two cannot be separated from the observed
means alone.

The practical consequence is narrow and firm. A threshold carried over from the
developmental cohort without re-estimating the intercept will misclassify in a
direction set by the sign of $\delta$, uniformly, for everyone. That is the
kind of error an external validation study is meant to find, and it is invisible
to any statistic computed within one cohort.
\end{reading}

Recalibration settles what the score means in haemophilia. Whether the score is
worth computing at all depends on what it predicts.


\subsection{Phase 4: Prognostic evaluation}

Once the translated measurement model has been established, the recalibrated $H_{\text{potential}}$ should undergo prospective prognostic evaluation. Disease-specific outcomes may include recurrent bleeding, progression of haemophilic arthropathy, unplanned hospitalisation, orthopaedic intervention, quality-of-life deterioration and all-cause mortality. The analytical strategy should mirror that employed in the present thesis: association, discrimination, calibration, comparative performance and clinical utility should all be evaluated using contemporary reporting standards. Maintaining methodological consistency across developmental and translational studies permits meaningful comparison of predictive performance while reducing analytical heterogeneity. Each of these terms has a counterpart in diagnostic test evaluation. Discrimination is the score's ability to rank people who go on to have the outcome above those who do not, and it is reported as the area under a receiver operating characteristic curve, read as for a diagnostic test. Calibration is agreement between the risks the score predicts and the rates observed. Two features of cohort follow-up also carry over. People who leave the registry before the end of follow-up are censored, so their outcome is unknown after they leave. And death is a competing risk for the other outcomes, since a person who dies can no longer be admitted or have a bleed, so the analysis must treat death as an event that ends observation of the outcome of interest and not as ordinary loss to follow-up.

The comparison that decides this phase is not whether the score predicts outcomes at all, since a score built partly from joint health and bleed rate can be expected to predict future bleeds. It is whether the score predicts them beyond what haemophilia care already measures. The reference model should therefore contain baseline factor level, joint score and inhibitor status, and the test is whether adding the score improves discrimination, calibration and net benefit over that reference. Two results would be informative in different ways. If the score adds nothing, the reading of Profile~B above applies: the framework would have been shown to add nothing to existing measures, which for clinical use is a verdict against it. If it adds something, the next question is where the gain comes from, and the systems reading of \S\ref{sec:hemo_systems} gives a specific expectation. The gain should come mainly from the social and systemic domains, and it should be larger in settings where supply is uncertain than in settings where it is secure. A gain spread evenly across settings, or concentrated in the biomedical domain, would suggest that the score is acting as a finer measure of severity and not as a measure of capacity in the sense this thesis intends.

Prognostic performance is necessary and not sufficient for use. Nothing in
the sequence may begin without permission, and none of it is worth much unless
it is reported so that others can check it.

\subsection{Phase 5: Governance, ethics and reporting}

Translation into clinical research requires governance structures extending beyond statistical validation alone. Ethical approval, informed consent, compliance with institutional review board requirements, secure management of participant information and adherence to applicable data-protection legislation must all be established before commencement of empirical studies. Transparent reporting is equally important, and future validation studies should be reported in accordance with internationally recognised guidance for prediction model research, including the TRIPOD Statement and subsequent updates \citep{collins_2015,moons_2019}. Such reporting facilitates reproducibility, critical appraisal and comparison with other prognostic models developed within haemophilia and related chronic diseases.

Although the present chapter focuses specifically on haemophilia, the proposed roadmap is intentionally generalisable. The same sequence of indicator mapping, expert consultation, psychometric evaluation, recalibration and prognostic validation could be applied to numerous chronic diseases, so that the framework establishes a general method for translating the $H_{\text{potential}}$, of which haemophilia is the first case, into other clinical domains while preserving its conceptual integrity. A dynamical specialisation of the framework to haemophilia, modelling clotting-factor kinetics, inhibitor development and gene therapy, is set out in Appendix~\ref{app:hemo_dynamics} as a separate future modelling direction. That work should be viewed as an extension of the current framework rather than a prerequisite for its validation: the present thesis is concerned with measurement and prognostic prediction, while mechanistic simulation provides an opportunity for future refinement once empirical validation has been completed.

What remains is to state plainly how much of the chapter rests on evidence and
how much on argument.

% =====================================================================
\section{Summary: What the Chapter Settles and What It Leaves Open}
\label{sec:hemo_close}
% =====================================================================

\noindent The chapter opened by asking whether a measure of health capacity built in a general ageing cohort can be carried into a disease it never saw and still mean the same thing. It has not answered that question with data, and it could not, since no haemophilia cohort is analysed in this thesis. What it has done is set out how the $H_{\text{potential}}$ would be translated into haemophilia while keeping its structure. Rather than building a new instrument for the disease, the strategy keeps the latent domains of Biomedical Potential, Social Modifiers and Systemic Burdens and replaces their observed indicators with variables that reflect the biology and health care of haemophilia.

\paragraph{What Is New in This Chapter.} Before this chapter, haemophilia care already had well-tested measures of its own, such as the Haemophilia Joint Health Score and the annualised bleed rate \citep{hilliard_2006}, and guidance on external validation already said that a prediction model moved to a new population must be validated and recalibrated there \citep{moons_2019,collins_2024}. The $H_{\text{potential}}$ itself had been specified and tested only in an ageing cohort. The chapter adds five things. First, a mapping of each domain onto indicators that haemophilia clinics already record, with a stated direction for each, set beside the general indicators they replace (Table~\ref{tab:hemo_mapping}, Table~\ref{tab:hemo-map}, Figure~\ref{fig:hemo-domains}), and with baseline factor level left out on purpose so that it can serve as the benchmark. Second, a formal statement of what replacing indicators requires (Assumption~\ref{ass:hem:substitution}) and a proof that, without shared anchors, a haemophilia cohort mean cannot be compared with an ageing cohort mean (Proposition~\ref{prop:hem:scalar}). Third, a consensus rule for the Delphi panel fixed before any rating is seen (Derivation~\ref{der:consensus}). Fourth, a reading of the haemophilia model of Appendix~\ref{app:hemo_dynamics} as the control loop the clinic runs, which shows that the tolerable inhibitor titre moves with the dose and that a binary inhibitor indicator therefore hides information the score needs. Fifth, a staged roadmap in which each stage has a stated product and a stated decision that can stop the programme (Figure~\ref{fig:hemo-roadmap}). The worked examples can be checked now: each is asserted by a named test in the verification suite (Appendix~\ref{app:verification}). The empirical claims cannot be checked until a haemophilia cohort exists; they are stated so that Phases~2 and~4 can refute them, and the conditions that would count as refutation are written down here, before the data.

Some matters are settled by argument alone. Derivation~\ref{der:indicator_substitution} and Assumption~\ref{ass:hem:substitution} state what indicator substitution requires and show that it buys correspondence of content, not a shared scale. Proposition~\ref{prop:hem:scalar} shows that, without scalar invariance or common anchors, a difference between a haemophilia cohort mean and an NHATS cohort mean carries no meaning as a difference in capability, which rules out a whole class of tempting claims before results arrive. Derivation~\ref{der:recalibration} and Worked Example~\ref{ex:hem:recal} separate the re-estimation of how capability is measured from the re-estimation of how a score maps to risk, and show that an intercept shift alone can move a cut-point enough to change who is flagged. Worked Example~\ref{ex:hem:mrb} shows that, under the realisation model, the funding setting alone can more than halve realised capability for a fixed biology. These results hold whatever the data later show, because they are properties of the model, not of the population.

Other matters rest on argument that the data could overturn. The case for the biological plausibility of each domain was made from what is known of haemophilia: Biomedical Potential through joint health, bleeding phenotype and inhibitor status; Social Modifiers through team-based care, adherence, caregiver support and health literacy; Systemic Burdens through treatment supply, health services and wider organisational factors that act whatever the biological severity. Plausibility is the weakest evidence the thesis accepts anywhere. It shows that the domains correspond to mechanisms; it does not show that the mechanisms combine as the realisation model says. The proposed indicator set likewise rests on a literature review until the Delphi consultation has tested it, and the consultation will establish agreement among clinicians at one centre, not measurement properties.

The systems reading of \S\ref{sec:hemo_systems} adds a third kind of result: a list of places to look. Read as a control system, haemophilia treatment has sharp thresholds observed through lagging measurements, an actuator that depends on supply and saturates at high titre, and a common cause, funding, that can move several domains at once. Each of these gives a specific expectation for the validation programme: that the score should move with supply before joint damage does, that domains will correlate more strongly where funding settings are mixed, and that a binary inhibitor indicator will weaken the biomedical domain. The same reading exposed a limit. Structural freedom at one time becomes biomedical potential at a later time, and a static score can only record the results of that transfer, not represent it.

What remains open is the empirical question itself. Whether the three domains reappear in a haemophilia cohort, whether the recalibrated score predicts bleeding, admission and joint deterioration beyond factor level and joint score, and whether any gain comes from the social and systemic domains as the framework predicts, can be answered only by the five phases of \S\ref{sec:hemo_roadmap}. Two design questions are also left open for decision before that work begins: whether people living with haemophilia should form a parallel Delphi panel, and how the registry will record the route by which each person reached care, so that the boundary of the cohort can be judged.

Success in haemophilia would be evidence that the $H_{\text{potential}}$ can be applied to a specific disease, and some evidence for its wider claim to measure latent health capacity across clinical settings, since the same sequence of steps could be applied to other chronic diseases. Failure would be informative too, and the chapter has stated in advance what each kind of failure would look like.

With this chapter each aim of the thesis has been addressed, Aim~4 by the translation set out here. The chapter does not show that the construct works in haemophilia; it shows what would have to be true for it to do so, and how that would be checked. The Discussion that follows in Chapter~\ref{ch:discussion} takes the thesis as a whole: it states what the findings would mean against the four aims, sets the $H_{\text{potential}}$ beside the measures it would have to sit with, and separates what the design earns from what it cannot.
